% =====================================================================
%  Benchmark Framework cho Streaming Clustering tren River
%  Bien dich:  pdflatex / xelatex / lualatex  (chay 2 lan de co muc luc)
%    pdflatex benchmark_framework_report.tex
%    pdflatex benchmark_framework_report.tex
% =====================================================================
\documentclass[11pt,a4paper]{article}

\usepackage{iftex}
\ifPDFTeX
  \usepackage[utf8]{inputenc}
  \usepackage[T5]{fontenc}
  \usepackage{lmodern}
\else
  \usepackage{fontspec}
  \setmainfont{Latin Modern Roman}
  \setmonofont{Latin Modern Mono}
\fi

\usepackage[margin=2.5cm]{geometry}
\usepackage{amsmath,amssymb,amsthm}
\usepackage{booktabs}
\usepackage{multirow}
\usepackage{array}
\usepackage{graphicx}
\usepackage{xcolor}
\usepackage{tikz}
\usepackage{pgfplots}
\pgfplotsset{compat=1.17}
\usetikzlibrary{arrows.meta,positioning,shapes.geometric,fit,backgrounds}
\usepackage{algorithm}
\usepackage{algpseudocode}
\usepackage{listings}
\usepackage{caption}
\usepackage[hidelinks]{hyperref}

\definecolor{codebg}{RGB}{246,247,249}
\definecolor{codekw}{RGB}{20,90,160}
\definecolor{codecm}{RGB}{110,120,130}
\lstset{
  basicstyle=\ttfamily\small,
  backgroundcolor=\color{codebg},
  keywordstyle=\color{codekw}\bfseries,
  commentstyle=\color{codecm}\itshape,
  frame=single,
  rulecolor=\color{codebg},
  breaklines=true,
  showstringspaces=false,
  columns=fullflexible,
  xleftmargin=6pt,
  framexleftmargin=6pt,
}

\newtheorem{proposition}{Mệnh đề}
\newtheorem{definition}{Định nghĩa}
\theoremstyle{remark}
\newtheorem{remark}{Nhận xét}

\newcommand{\R}{\mathbb{R}}
\newcommand{\E}{\mathbb{E}}
\newcommand{\xt}{\mathbf{x}}
\newcommand{\ind}[1]{\mathbf{1}\!\left[#1\right]}

\title{\bfseries Một framework benchmark cho \texttt{river.cluster}:\\
Giao thức đánh giá, độ đo, và những cái bẫy đo lường}
\author{Báo cáo kỹ thuật --- \texttt{river-customized-benchmark-framework}}
\date{\today}

\begin{document}
\maketitle

\begin{abstract}
\noindent
Báo cáo này mô tả một \textbf{framework đo lường} cho mô-đun
\texttt{river.cluster}: nó đo cái gì, đo bằng công thức nào, chạy theo giao thức
nào, và tại sao một số cách đo tưởng như hiển nhiên lại cho kết luận sai. Sản
phẩm ở đây là \emph{bộ máy đo}, không phải một bảng xếp hạng thuật toán: việc
chọn và tối ưu siêu tham số cũng như cấu hình dataset thuộc về người dùng, còn
framework có nhiệm vụ đảm bảo mọi cấu hình đều được chạy dưới cùng một giao thức
và được báo cáo trung thực.

Chúng tôi hình thức hoá giao thức \emph{test-then-train}; phân biệt rõ hai họ chỉ
số chất lượng --- \emph{luỹ kế} và \emph{theo cửa sổ} --- và chứng minh rằng lấy
trung bình đường luỹ kế không phải là ``trung bình theo thời gian''
(Mệnh đề~\ref{prop:cumbias}). Đóng góp chính về phương pháp là \textbf{chẩn đoán
cửa sổ thoái hoá}: trên một cửa sổ chỉ chứa một lớp ground truth, ARI suy biến
thành hàm nhị phân $\{0,1\}$ chỉ phụ thuộc vào việc mô hình có gộp toàn bộ cửa sổ
lại hay không, hoàn toàn độc lập với chất lượng phân cụm
(Mệnh đề~\ref{prop:degenerate}). Trên KDDCup99, $72\%$ số cửa sổ kích thước
$1000$ rơi vào trường hợp này, khiến bảng xếp hạng theo hai chỉ số gần như phản
bác nhau ($\rho_s = -0.8$); framework do đó đo và cảnh báo trực tiếp tình trạng
này thay vì để nó lọt qua thành một con số trông như kết quả.

Framework bao phủ \textbf{cả bảy} clusterer của River, kể cả hai cái phá vỡ giao
thức thông thường (ODAC gom cụm \emph{các biến}; TextClust ăn \emph{token
counts}), và toàn bộ 25 bước tiền xử lý không giám sát của River, ghép chuỗi
được. Trong quá trình xây dựng, hai lỗi của River được phát hiện và \textbf{chỉ
được ghi lại} --- cây nguồn River không bị chỉnh sửa --- kèm phương án né phía
framework (Mục~\ref{sec:riverbugs}).
\end{abstract}

\tableofcontents
\newpage

% =====================================================================
\section{Giới thiệu}
% =====================================================================

Phân cụm trên luồng dữ liệu (\emph{streaming clustering}) khác phân cụm truyền
thống ở ba điểm khiến việc đánh giá trở nên khó: dữ liệu đến tuần tự và chỉ được
xem một lần, mô hình phải trả lời ngay tại thời điểm nhận mẫu, và tài nguyên
(thời gian, bộ nhớ) là ràng buộc cứng chứ không phải yếu tố phụ. Vì vậy một
benchmark cho lớp thuật toán này không thể chỉ báo một con số chất lượng cuối
cùng; nó phải mô tả được \emph{quỹ đạo} của cả chất lượng lẫn tài nguyên dọc
theo luồng.

Framework được mô tả ở đây trả lời câu hỏi:

\begin{quote}
\itshape
Thuật toán nào cho phân cụm tốt hơn, xử lý nhanh hơn, tốn ít bộ nhớ hơn, và
mở rộng tốt hơn trên dữ liệu luồng?
\end{quote}

Nó gồm bốn nhóm chỉ số --- chất lượng, hiệu năng tính toán, độ phức tạp mô
hình, khả năng mở rộng --- cộng một nhóm thứ năm về độ ổn định qua nhiều seed.
Phần đóng góp mà chúng tôi cho là quan trọng nhất lại không nằm ở việc thêm chỉ
số, mà ở việc chỉ ra \emph{khi nào một chỉ số mất ý nghĩa} và đo trực tiếp điều
đó (Mục~\ref{sec:degenerate}).

\paragraph{Bố cục.}
Mục~\ref{sec:notation} thiết lập ký hiệu và giao thức trực tuyến.
Mục~\ref{sec:arch} mô tả kiến trúc.
Mục~\ref{sec:scaler} trình bày các bộ chuẩn hoá trực tuyến kèm công thức đệ quy.
Mục~\ref{sec:algos} tóm tắt bốn thuật toán được benchmark.
Mục~\ref{sec:protocol} là phần trọng tâm: giao thức đánh giá và toàn bộ công
thức độ đo.
Mục~\ref{sec:degenerate} trình bày và chứng minh vấn đề cửa sổ thoái hoá.
Mục~\ref{sec:pseudocode} cho mã giả đầy đủ.
Mục~\ref{sec:results} là kết quả thực nghiệm.
Mục~\ref{sec:usage} hướng dẫn chạy.
Mục~\ref{sec:limits} nêu hạn chế.
Hình~\ref{fig:arch} tóm tắt kiến trúc và Hình~\ref{fig:timeline} tóm tắt hai
trục thời gian của phép đánh giá; Mục~\ref{sec:fullbench} là kết quả tổng hợp
trên cả bốn thuật toán.

% =====================================================================
\section{Ký hiệu và mô hình bài toán}
\label{sec:notation}
% =====================================================================

\subsection{Luồng dữ liệu}

Một \emph{luồng} là dãy hữu hạn có thứ tự
\begin{equation}
  \mathcal{S} \;=\; \big( (\xt_t,\, y_t) \big)_{t=1}^{T},
  \qquad
  \xt_t \in \R^{d},
  \qquad
  y_t \in \mathcal{Y} \cup \{\bot\},
\end{equation}
trong đó $\xt_t$ là véc-tơ đặc trưng tại bước $t$ (trong River được biểu diễn
bằng \texttt{dict} ánh xạ tên đặc trưng sang giá trị), $y_t$ là nhãn ground truth
dùng \emph{duy nhất} để chấm điểm, và $\bot$ ký hiệu mẫu không có nhãn (ví dụ
nhiễu nền do bộ sinh dữ liệu chèn vào). Thứ tự $t$ là thứ tự bản ghi trong tập
dữ liệu --- không xáo trộn, vì thứ tự đến \emph{chính là} một phần của bài toán.

\subsection{Bộ phân cụm trực tuyến}

Một bộ phân cụm trực tuyến là cặp ánh xạ hoạt động trên trạng thái $\theta$:
\begin{align}
  \texttt{predict}&:\ (\theta,\ \xt) \;\mapsto\; \hat{y} \in \mathbb{N}, \\
  \texttt{learn}&:\ (\theta,\ \xt) \;\mapsto\; \theta'.
\end{align}
Khác với học có giám sát, \texttt{learn} \emph{không} nhận $y$. Nhãn cụm $\hat y$
chỉ là định danh, không mang ngữ nghĩa: hai phân hoạch chỉ khác nhau ở cách đánh
số cụm phải được coi là giống hệt nhau. Đây chính là lý do mọi độ đo ngoại
(external metric) ở Mục~\ref{sec:external} đều bất biến với hoán vị nhãn.

\subsection{Giao thức test-then-train}

Tại mỗi bước $t$, framework thực hiện đúng ba thao tác theo thứ tự:
\begin{equation}
  \label{eq:ttt}
  \boxed{\;
  \xt_t \;\xrightarrow{\ \text{scale}\ }\; \tilde\xt_t
  \;\xrightarrow{\ \texttt{predict}\ }\; \hat y_t
  \;\xrightarrow{\ \texttt{learn}\ }\; \theta_t
  \;}
\end{equation}
Dự đoán xảy ra \emph{trước} khi học, đúng như một hệ thống chạy thật: khi mẫu
đến, ta phải trả lời ngay, chưa kịp học từ chính nó. Cách đánh giá này gọi là
\emph{prequential} (predictive sequential): mỗi mẫu vừa là dữ liệu kiểm thử vừa
là dữ liệu huấn luyện, và không cần tách tập test riêng.

\begin{figure}[htbp]
\centering
\begin{tikzpicture}[
  node distance=3mm,
  box/.style={draw,rounded corners=2pt,minimum height=7mm,minimum width=15mm,
              font=\small,fill=blue!5},
  lab/.style={font=\scriptsize\itshape,text=gray!70!black},
]
% timeline
\draw[-{Stealth[length=2mm]},thick] (0,0) -- (14.6,0);
\node[lab] at (14.9,-0.32) {$t$};

% window blocks
\foreach \i/\x in {1/0.2, 2/3.8, 3/7.4, 4/11.0} {
  \draw[fill=blue!8,draw=blue!40] (\x,0.15) rectangle (\x+3.4,1.05);
  \node[font=\scriptsize] at (\x+1.7,0.6) {cửa sổ $W_{\i}$ ($m$ mẫu)};
}
% flush markers
\foreach \x in {3.6, 7.2, 10.8, 14.4} {
  \draw[dashed,red!70] (\x,-0.55) -- (\x,1.25);
}
\node[lab,text=red!70,anchor=west] at (3.65,1.45) {flush: chốt $M^{\mathrm{win}}_k$, reset metric cửa sổ};

% checkpoints
\foreach \x in {1.9, 5.5, 9.1, 12.7} {
  \draw[-{Stealth[length=1.5mm]},orange!85!black,thick] (\x,-1.15) -- (\x,-0.62);
  \node[circle,fill=orange!85!black,inner sep=1pt] at (\x,-0.55) {};
}
\node[lab,text=orange!70!black,anchor=west] at (0.0,-1.45)
  {checkpoint: đo runtime luỹ kế, bộ nhớ mô hình, số micro-cluster};

% cumulative bar
\draw[{Stealth[length=2mm]}-{Stealth[length=2mm]},thick,green!45!black]
  (0.2,1.85) -- (14.4,1.85);
\node[lab,text=green!35!black,anchor=south] at (7.3,1.9)
  {metric luỹ kế $M^{\mathrm{cum}}$: không bao giờ reset, chạy suốt luồng};
\end{tikzpicture}
\caption{Hai trục thời gian của phép đánh giá. Các metric \emph{luỹ kế} chạy
liên tục từ đầu luồng; các metric \emph{theo cửa sổ} được reset ở mỗi biên. Các
checkpoint đo tài nguyên nằm trên một lưới độc lập, thưa hơn, vì phép đo bộ nhớ
rất đắt (Mục~\ref{sec:efficiency}).}
\label{fig:timeline}
\end{figure}

% =====================================================================
\section{Kiến trúc framework}
\label{sec:arch}
% =====================================================================

\begin{figure}[htbp]
\centering
\begin{tikzpicture}[
  node distance=8mm and 12mm,
  mod/.style={draw,rounded corners=3pt,minimum height=11mm,minimum width=27mm,
              align=center,font=\small,fill=blue!6,draw=blue!45},
  out/.style={draw,rounded corners=3pt,minimum height=11mm,minimum width=27mm,
              align=center,font=\small,fill=green!7,draw=green!45!black},
  ar/.style={-{Stealth[length=2.2mm]},thick,gray!60!black},
]
\node[mod] (ds)  {\texttt{datasets/}\\[-1pt]\scriptsize luồng $(\xt_t,y_t)$};
\node[mod,right=of ds] (set) {\texttt{Setting}\\[-1pt]\scriptsize thuật toán + params\\[-2pt]\scriptsize + scaler + scaler\_params};
\node[mod,right=of set] (ad) {\texttt{ClustererAdapter}\\[-1pt]\scriptsize scale $\to$ predict $\to$ learn};
\node[mod,below=14mm of ad] (run) {\texttt{run\_once}\\[-1pt]\scriptsize runner trực tuyến};
\node[out,left=of run] (ev) {\texttt{evaluation/}\\[-1pt]\scriptsize metric + checkpoint};
\node[out,left=of ev] (rep) {\texttt{reports/}\\[-1pt]\scriptsize bảng, đồ thị, JSON};

\draw[ar] (ds) -- (set);
\draw[ar] (set) -- (ad);
\draw[ar] (ad) -- (run);
\draw[ar] (run) -- (ev);
\draw[ar] (ev) -- (rep);
\draw[ar] (ds.south) |- ([yshift=-3mm]run.east) -- (run.east);
\end{tikzpicture}
\caption{Luồng dữ liệu qua framework. Mọi \texttt{Setting} tiêu thụ \emph{cùng
một} đối tượng dataset, và dataset trả về iterator mới ở mỗi lần lặp, nên mọi
thuật toán thấy đúng cùng một dãy bản ghi theo đúng thứ tự.}
\label{fig:arch}
\end{figure}

Bảng~\ref{tab:modules} liệt kê vai trò từng mô-đun.

\begin{table}[htbp]
\centering
\small
\caption{Các mô-đun của framework.}
\label{tab:modules}
\begin{tabular}{@{}lp{10.2cm}@{}}
\toprule
\textbf{Mô-đun} & \textbf{Trách nhiệm} \\
\midrule
\texttt{datasets/} & Nguồn luồng \emph{lặp lại được}: sinh Gaussian tổng hợp, KDDCup99 (34 đặc trưng liên tục), KDDCup98 (56 trường RFM). \\
\texttt{algorithms/registry} & Đăng ký thuật toán + preset tham số; nhà máy tạo scaler theo tên kèm tham số. \\
\texttt{algorithms/adapter} & Bọc mọi \texttt{river.base.Clusterer} sau một giao diện chung; gắn scaler; trích số micro/macro-cluster và bộ nhớ. \\
\texttt{evaluation/runner} & Vòng lặp trực tuyến; quản lý cửa sổ đánh giá và checkpoint tài nguyên. \\
\texttt{evaluation/clustering\_quality} & Metric streaming của River + hàm tham chiếu theo lô để đối chứng. \\
\texttt{experiments/harness} & Chạy nhiều setting trên cùng luồng; lặp seed; tổng hợp mean $\pm$ std. \\
\texttt{reports/} & Bảng markdown, đồ thị matplotlib, bản ghi môi trường chạy. \\
\bottomrule
\end{tabular}
\end{table}

% =====================================================================
\section{Tiền xử lý trực tuyến}
\label{sec:scaler}
% =====================================================================

Bộ chuẩn hoá (scaler) đứng ngay trước bộ phân cụm trong~\eqref{eq:ttt}. Với các
thuật toán dựa trên khoảng cách, nó quyết định thang đo mà mọi ngưỡng bán kính
($\varepsilon$ của DenStream, \texttt{clustering\_threshold} của DBSTREAM) được
diễn giải --- nói cách khác, scaler và tham số bán kính \emph{không độc lập}.

\subsection{Thứ tự học rồi biến đổi}

Với mỗi mẫu, adapter thực hiện
\begin{equation}
  \label{eq:scaleorder}
  \theta^{\text{sc}}_t \leftarrow \texttt{learn}(\theta^{\text{sc}}_{t-1}, \xt_t),
  \qquad
  \tilde\xt_t \leftarrow \texttt{transform}(\theta^{\text{sc}}_t, \xt_t).
\end{equation}
Tức là mẫu hiện tại \emph{có} đóng góp vào thống kê trước khi chính nó được biến
đổi. Thứ tự này trùng khớp với \texttt{river.compose.Pipeline}, nơi các bước
biến đổi không giám sát học trước rồi mới biến đổi; do đó adapter tương đương với
\texttt{StandardScaler() | cluster.X()}.

\begin{remark}[Khởi động lạnh]
Ở $t=1$ ta có $\sigma^2=0$ nên mọi đặc trưng đều được ánh xạ về $0$. Vài trăm mẫu
đầu vì thế được chuẩn hoá bằng thống kê chưa ổn định. Đây là đặc tính cố hữu của
chuẩn hoá trực tuyến, không phải lỗi; nhưng nó có nghĩa là mọi chỉ số luỹ kế đều
mang theo giai đoạn khởi động lạnh này mãi mãi (xem Mục~\ref{sec:cumwin}).
\end{remark}

\subsection{StandardScaler: đệ quy Welford}

Với mỗi đặc trưng $j$, gọi $n$ là số quan sát đã thấy. River cập nhật theo
đệ quy Welford:
\begin{align}
  n_t &= n_{t-1} + 1, \\
  \mu_t &= \mu_{t-1} + \frac{x_t - \mu_{t-1}}{n_t},
  \label{eq:welford-mean}\\
  \sigma^2_t &= \sigma^2_{t-1}
    + \frac{(x_t - \mu_{t-1})(x_t - \mu_t) - \sigma^2_{t-1}}{n_t},
  \label{eq:welford-var}
\end{align}
cho phương sai tổng thể (\texttt{ddof}$=0$). Đệ quy~\eqref{eq:welford-mean}
và~\eqref{eq:welford-var} chỉ cần $\mathcal{O}(1)$ bộ nhớ cho mỗi đặc trưng và
không bao giờ phải lưu lại quan sát nào. Phép biến đổi là
\begin{equation}
  \label{eq:zscore}
  z_t \;=\;
  \begin{cases}
    \dfrac{x_t - \mu_t}{\sqrt{\sigma^2_t}}, & \sigma^2_t > 0,\\[6pt]
    0, & \text{ngược lại}.
  \end{cases}
\end{equation}

\begin{remark}[Điều kiện $\sigma^2_t>0$ chứ không phải $\sigma^2_t \neq 0$]
\label{rem:guard}
Trong bản River gốc, điều kiện trong~\eqref{eq:zscore} được viết dưới dạng kiểm
tra ``khác không''. Với
một đặc trưng hằng số, sai số dấu phẩy động có thể đẩy $\sigma^2$ xuống một giá
trị \emph{âm} rất nhỏ; khi đó $(\sigma^2)^{1/2}$ trong Python trả về \emph{số
phức}, và giá trị phức này lọt vào mô hình rồi mới gây lỗi ở tận hàm tính khoảng
cách. Chúng tôi đo được hiện tượng này trên KDDCup99; River \emph{không} được sửa,
lỗi chỉ được ghi lại và né ở phía framework (Mục~\ref{sec:bug1}).
\end{remark}

\subsection{Biến thể cửa sổ trượt}

Khi đặt \texttt{window\_size} $=W$, trung bình và phương sai được tính trên $W$
quan sát gần nhất:
\begin{equation}
  \mu^{(W)}_t = \frac{1}{|B_t|}\sum_{s \in B_t} x_s,
  \qquad
  \big(\sigma^{(W)}_t\big)^2 = \frac{1}{|B_t|}\sum_{s \in B_t} (x_s - \mu^{(W)}_t)^2,
  \qquad
  B_t = \{\,t-W+1,\dots,t\,\} \cap \mathbb{N}^{+}.
\end{equation}
Điều này đòi hỏi \emph{lưu trữ} $W$ quan sát cho \emph{mỗi} đặc trưng, tức
$\mathcal{O}(dW)$ bộ nhớ --- một cái giá rất thật mà chúng tôi đo được ở
Mục~\ref{sec:ablation}.

\subsection{AdaptiveStandardScaler: quên theo hàm mũ}

Đây là bộ chuẩn hoá duy nhất \emph{quên} dữ liệu cũ mà chỉ tốn $\mathcal{O}(d)$
bộ nhớ. Với hệ số phai $\alpha \in [0,1]$, đặt $\delta_t = x_t - m_{t-1}$:
\begin{align}
  m_t &= m_{t-1} + \alpha\,\delta_t, & m_1 &= x_1,
  \label{eq:ewmean}\\
  v_t &= (1-\alpha)\big(v_{t-1} + \alpha\,\delta_t^2\big), & v_1 &= 0,
  \label{eq:ewvar}
\end{align}
rồi $z_t = (x_t - m_t)/\sqrt{v_t}$. Các công thức~\eqref{eq:ewmean}
và~\eqref{eq:ewvar} đã được chúng tôi kiểm chứng trùng khớp từng chữ số với
\texttt{river.stats.EWMean} và \texttt{river.stats.EWVar}. Trọng số của quan sát
cách hiện tại $k$ bước suy giảm như $\alpha(1-\alpha)^{k}$, cho ``bán kính nhớ''
hiệu dụng cỡ $1/\alpha$ mẫu.

\subsection{Các scaler còn lại}

\begin{align}
  \text{MinMax:}\quad & z_t = \frac{x_t - \min_{s\le t} x_s}{\max_{s\le t} x_s - \min_{s\le t} x_s},\\
  \text{MaxAbs:}\quad & z_t = \frac{x_t}{\max_{s \le t} |x_s|},\\
  \text{Robust:}\quad & z_t = \frac{x_t - Q_{0.5}}{Q_{q_{\sup}} - Q_{q_{\inf}}},
    \quad (q_{\inf},q_{\sup}) = (0.25,\,0.75)\ \text{mặc định},\\
  \text{Normalizer:}\quad & \tilde\xt_t = \frac{\xt_t}{\|\xt_t\|_p} \quad (\text{không có trạng thái}).
\end{align}

Bảng~\ref{tab:scalers} tóm tắt các scaler theo hai tiêu chí quyết định: khả năng
quên và chi phí bộ nhớ.

\begin{table}[htbp]
\centering
\small
\caption{Các scaler trực tuyến: điều khác biệt thực chất là \emph{nhớ bao nhiêu}
và \emph{tốn bao nhiêu bộ nhớ}. $d$ là số đặc trưng, $W$ là kích thước cửa sổ.}
\label{tab:scalers}
\begin{tabular}{@{}llcc@{}}
\toprule
\textbf{Scaler} & \textbf{Thống kê} & \textbf{Quên dữ liệu cũ?} & \textbf{Bộ nhớ} \\
\midrule
\texttt{standard} & mean/var Welford & không & $\mathcal{O}(d)$ \\
\texttt{standard(window\_size=$W$)} & mean/var trượt & có, cứng & $\mathcal{O}(dW)$ \\
\texttt{minmax}, \texttt{maxabs} & min/max (hoặc trượt) & không (trừ khi trượt) & $\mathcal{O}(d)$ / $\mathcal{O}(dW)$ \\
\texttt{robust} & trung vị + IQR & không & $\mathcal{O}(d)$ xấp xỉ \\
\texttt{adaptive\_standard} & EW mean/var & \textbf{có}, mềm & $\mathcal{O}(d)$ \\
\texttt{normalizer} & không có trạng thái & --- & $\mathcal{O}(1)$ \\
\bottomrule
\end{tabular}
\end{table}

\subsection{Không chỉ scaler: toàn bộ bước tiền xử lý không giám sát}
\label{sec:preproc-all}

Slot đứng trước bộ phân cụm nhận \emph{mọi} transformer không giám sát của River,
tổng cộng 25 bước. Benchmark clustering không có $y$ lúc học nên mọi thứ cần
target bị loại theo định nghĩa; phần còn lại lấp những lỗ hổng thật, và bỏ chúng
đi thì có những câu hỏi không đặt ra được (Bảng~\ref{tab:preproc}). Nhận
xét~\ref{rem:onehot} là ví dụ cụ thể về việc chọn nhầm bước trong nhóm này phá
hỏng phép đo ra sao.

\begin{table}[htbp]
\centering
\small
\caption{Các nhóm bước tiền xử lý và lỗ hổng chúng lấp. Cố ý vắng mặt:
\texttt{SelectKBest}, \texttt{PoissonInclusion}, các \texttt{Target*} (cần
target); \texttt{PredClipper} (bọc regressor); \texttt{PolynomialExtender} (nhân
số chiều và bóp méo đúng cái khoảng cách mà bộ phân cụm dựa vào).}
\label{tab:preproc}
\begin{tabular}{@{}lp{8.4cm}@{}}
\toprule
\textbf{Nhóm} & \textbf{Vì sao benchmark clustering cần} \\
\midrule
Scaling & mọi tham số bán kính ($\varepsilon$ của DenStream,
\texttt{clustering\_threshold} của DBSTREAM) được đọc theo đơn vị scaler tạo ra \\
Missing values & stream thật có thiếu giá trị, và quy ước của River là thiếu
$=$ \emph{vắng key}, nên bộ phân cụm sẽ so hai điểm trên hai tập key khác nhau \\
Categorical & không có thì phải \emph{vứt} cột nominal --- KDDCup99 mất 3 cột,
CoverType mất 44 \\
Projection & giảm chiều bảo toàn khoảng cách (Johnson--Lindenstrauss), công cụ
tiêu chuẩn cho clustering luồng nhiều chiều \\
Variance filter & cột hằng số không đóng góp gì vào khoảng cách và còn phá
scaler --- chính một cột hằng số đẩy rolling \texttt{StandardScaler} của River
xuống phương sai âm (Mục~\ref{sec:riverbugs}) \\
Text & TextClust ăn token counts; \texttt{lda} biến text thành vector dày để các
bộ phân cụm \emph{số} cũng xử lý được text \\
Kernel & random Fourier features cho phép phương pháp tâm-cụm tách được cụm
không có dạng cầu \\
Routing & chọn/bỏ cột trước một bước không tự phân biệt được kiểu dữ liệu \\
\bottomrule
\end{tabular}
\end{table}

\paragraph{Ghép chuỗi.}
Pipeline hữu ích hiếm khi chỉ một bước: dữ liệu thô thường cần điền thiếu và mã
hoá phạm trù \emph{trước khi} chuẩn hoá có ý nghĩa, còn luồng nhiều chiều muốn
bỏ cột hằng số rồi chiếu xuống trước khi tính bất kỳ khoảng cách nào. Các bước
ghép bằng \texttt{|}, chạy trái sang phải:

\begin{lstlisting}[language=Python]
make_preprocessor("impute|one_hot_categorical|standard")
make_preprocessor("variance_threshold|sparse_projection|standard",
                  sparse_projection={"n_components": 16})
\end{lstlisting}

Một bước nhận tham số phẳng; một chuỗi bắt buộc lồng tham số theo tên bước, vì
hai bước trong cùng chuỗi rất dễ trùng tên tham số (\texttt{seed}) và không có
cách nào đoán đúng chủ sở hữu.

\begin{remark}[Bẫy \texttt{one\_hot}]
\label{rem:onehot}
\texttt{OneHotEncoder} của River mã hoá \emph{mọi} key nó nhận. Trên một bản ghi
lẫn số và chữ, nó biến cột số thành một indicator cho \emph{mỗi giá trị từng
thấy}: với \texttt{duration} của KDDCup99 là hàng chục nghìn cột, mỗi cột mang
đúng một quan sát, và mọi khoảng cách sụp về gần như nhau. Framework do đó cung
cấp \texttt{one\_hot\_categorical}, chỉ mã hoá giá trị không phải số:
\[
\begin{array}{ll}
\texttt{one\_hot} & \to \{\texttt{dur\_1.0}: 0,\ \texttt{dur\_5.0}: 0,\ \texttt{dur\_7.0}: 1,\ \dots\}\\
\texttt{one\_hot\_categorical} & \to \{\texttt{dur}: 7.0,\ \texttt{proto\_tcp}: 0,\ \texttt{proto\_udp}: 1\}
\end{array}
\]
\end{remark}

% =====================================================================
\section{Các thuật toán được benchmark}
\label{sec:algos}
% =====================================================================

Framework bao phủ \textbf{cả bảy} clusterer trong \texttt{river.cluster}
(Bảng~\ref{tab:protocols}), với các bước tiền xử lý ở Mục~\ref{sec:preproc-all}.
Điều này không phải chuyện đếm số: hai trong bảy phá vỡ giao thức ``gán mẫu này
vào một cụm'' mà phần lớn framework benchmark ngầm giả định, và nếu giả định đó không
được kiểm tra thì hoặc là chương trình sập, hoặc --- tệ hơn --- một điểm chất
lượng vô nghĩa được ghi ra như thể là kết quả.

\begin{table}[htbp]
\centering
\small
\caption{Ba dạng giao thức trong \texttt{river.cluster} và cách framework xử lý.}
\label{tab:protocols}
\begin{tabular}{@{}llll@{}}
\toprule
\textbf{Thuật toán} & \textbf{Gán mẫu?} & \textbf{Đầu vào} & \textbf{Được chấm bằng} \\
\midrule
CluStream    & có & vector số & ARI/NMI/AMI/purity + tài nguyên \\
DenStream    & có & vector số & ARI/NMI/AMI/purity + tài nguyên \\
DBSTREAM     & có & vector số & ARI/NMI/AMI/purity + tài nguyên \\
STREAMKMeans & có & vector số & ARI/NMI/AMI/purity + tài nguyên \\
KMeans       & có & vector số & ARI/NMI/AMI/purity + tài nguyên \\
ODAC         & \textbf{không} & vector số & \textbf{cấu trúc} + tài nguyên \\
TextClust    & có & \textbf{token counts} & ARI/NMI/AMI/purity + tài nguyên \\
\bottomrule
\end{tabular}
\end{table}

\paragraph{ODAC: gom cụm các \emph{biến}.}
ODAC dựng cây phân cấp trên \emph{tập đặc trưng}, không phải trên các mẫu; hàm
\texttt{predict\_one} của nó \texttt{raise NotImplementedError} theo thiết kế.
Không có nhãn per-sample thì ARI/NMI/AMI/purity \emph{không định nghĩa được}, nên
framework bỏ nó khỏi mọi bảng chất lượng --- đó là đúng, không phải lỗi --- và
chấm nó bằng \emph{tóm tắt cấu trúc}: số cụm biến, số cụm còn hoạt động, chiều
cao cây. Cột ``Macro $K$'' cũng để trống cho ODAC, vì \texttt{n\_clusters} của nó
đếm nhóm \emph{biến}: khác đơn vị, đặt cạnh số cụm mẫu của các thuật toán khác là
mời người đọc so hai thứ không so được.

\begin{remark}[River là thẩm quyền về giao thức, không phải siêu dữ liệu của ta]
Nếu registry khai sai rằng một mô hình phân hoạch mẫu, mô hình đó vẫn phát hiện
\texttt{NotImplementedError} giữa chừng và tự hạ cờ. Điều này quan trọng vì các
đối tượng metric khi đó chưa nhận cập nhật nào mà vẫn trả về giá trị: ARI trả
$1.0$ trên bảng liên hợp rỗng theo quy ước $0/0$. Nghĩa là một mô hình
\emph{chưa hề dự đoán gì} sẽ được ghi điểm tuyệt đối. Runner do đó đọc lại cờ ở
cuối lượt chạy và huỷ bỏ các metric thay vì báo cáo chúng.
\end{remark}

\paragraph{TextClust: đầu vào không phải toạ độ.}
TextClust nhận dictionary tần suất từ và so sánh bằng cosine trên tf--idf, không
phải một metric trên toạ độ. Nó cần bước \texttt{bag\_of\_words} và một text
stream; chuẩn hoá tần suất từ bằng \texttt{StandardScaler} là vô nghĩa. Nó cũng
nhận vị trí stream qua tham số \texttt{t} khi bật chế độ fading theo thời gian
thực --- thiếu \texttt{t} thì River báo
\texttt{TypeError: unsupported operand type(s) for -}, rất xa nguyên nhân. Adapter
quyết định một lần bằng nội quan chữ ký hàm rồi truyền vào, thay vì đoán ở từng
mẫu.

\paragraph{Mô hình hai pha.}
Năm thuật toán trên vector số đều theo mô hình \emph{hai pha}: pha trực tuyến duy
trì một tóm tắt nhẹ của luồng (các \emph{micro-cluster}), pha ngoại tuyến gộp các
tóm tắt đó thành cụm cuối (\emph{macro-cluster}) khi được hỏi.

\paragraph{CluStream.}
Duy trì tối đa $q$ micro-cluster, mỗi cái là véc-tơ đặc trưng cụm
\begin{equation}
  \mathrm{CF} = \big(\overline{CF2^x},\ \overline{CF1^x},\ CF2^t,\ CF1^t,\ n\big),
\end{equation}
với $\overline{CF1^x}=\sum \xt_i$, $\overline{CF2^x}=\sum \xt_i^2$ (theo từng
toạ độ). Nhờ tính cộng tính của $\mathrm{CF}$, tâm và bán kính tính được trong
thời gian hằng:
\begin{equation}
  \mathbf{c} = \frac{\overline{CF1^x}}{n},
  \qquad
  r = \sqrt{\frac{\overline{CF2^x}}{n} - \left(\frac{\overline{CF1^x}}{n}\right)^{\!2}}.
\end{equation}

\paragraph{DenStream.}
Dựa trên mật độ, với suy giảm theo thời gian. Trọng số của một mẫu tuổi $\Delta t$
là $f(\Delta t) = 2^{-\lambda \Delta t}$ với $\lambda$ là \texttt{decaying\_factor}.
Một micro-cluster có trọng số $w$ được coi là \emph{lõi} nếu $w \ge \beta\mu$.
Một mẫu được thu nhận vào micro-cluster gần nhất khi
\begin{equation}
  \label{eq:denstream-eps}
  r\big(\text{mc} \cup \{\xt\}\big) \;\le\; \varepsilon ,
\end{equation}
và hai cụm lõi được gộp khi khoảng cách tâm nhỏ hơn $2\varepsilon$.

\begin{remark}[$\varepsilon$ phải khớp với thang sau tiền xử lý]
\label{rem:eps}
Sau \texttt{StandardScaler}, mỗi toạ độ có phương sai xấp xỉ $1$, nên khoảng
cách Euclid điển hình giữa hai điểm trong $\R^{d}$ là
$\E\|\xt-\xt'\|_2 \approx \sqrt{2d}$. Với KDDCup99 ($d=34$) con số này là
$\sqrt{68}\approx 8.2$: một $\varepsilon$ lớn hơn ngưỡng đó khiến
điều~\eqref{eq:denstream-eps} hầu như luôn đúng và toàn bộ luồng sụp về một
micro-cluster. Đây là một \emph{quan hệ} mà framework phơi bày qua cột
$\max|\mathcal{M}|$, chứ không phải một giá trị framework tự chọn: hiệu chuẩn
$\varepsilon$ là việc của người dùng, trên luồng hiệu chuẩn riêng.
\end{remark}

\paragraph{DBSTREAM.}
Cũng dựa trên mật độ nhưng dùng \emph{mật độ chia sẻ} giữa các micro-cluster.
Mẫu được gán vào mọi micro-cluster nằm trong bán kính
\texttt{clustering\_threshold}; trọng số suy giảm theo \texttt{fading\_factor}.
Tâm được cập nhật theo quy tắc kéo $\mathbf{c} \leftarrow \mathbf{c} + \eta(\xt - \mathbf{c})$.

\paragraph{STREAMKMeans và KMeans.}
STREAMKMeans tích luỹ luồng thành các khối kích thước $B$; mỗi khi đầy một khối
nó chạy một \texttt{KMeans} trên khối rồi nạp các tâm thu được vào một
\texttt{KMeans} chạy liên tục. Vì thế mô hình chưa có tâm nào cho tới khi khối
đầu tiên hoàn tất --- giai đoạn khởi động lạnh này được framework đếm và báo cáo
tường minh ở cột \emph{Cold-start}, để nó không âm thầm bóp méo các metric.

% =====================================================================
\section{Nguồn dữ liệu}
\label{sec:datasets}
% =====================================================================

Một dataset trong framework là nguồn \emph{lặp lại được} của các cặp $(\xt, y)$.
Tính lặp lại được chính là thứ cho phép mọi setting tiêu thụ đúng cùng một luồng:
đối tượng dataset trả về một iterator mới ở mỗi lần lặp. Nhãn $y$ chỉ dùng để
chấm điểm, và $y = \bot$ nghĩa là ``đừng chấm mẫu này''.

\subsection{Bộ sinh tổng hợp}

Theo cấu trúc của \texttt{RandomRBFGeneratorEvents} trong MOA. Điểm được sinh
dưới dạng \emph{hướng ngẫu nhiên $\times$ biên độ Gaussian} chứ không phải
Gaussian đẳng hướng, và khác biệt này có ý nghĩa: nó giữ kích thước vật lý của
cụm \textbf{độc lập với} $d$, nên phép quét ``chất lượng theo số chiều'' chỉ thay
đổi đúng số chiều thay vì âm thầm thổi mọi cụm lên $\sqrt{d}$ lần.

Ground truth là \textbf{id của cụm sinh ra điểm}, không phải một class label gộp
nhiều cụm --- với bài toán clustering thì khôi phục đúng thành phần hỗn hợp mới
là mục tiêu.

\begin{table}[htbp]
\centering
\small
\caption{Các núm điều khiển của bộ sinh.}
\label{tab:generator}
\begin{tabular}{@{}lp{8.6cm}@{}}
\toprule
\textbf{Tham số} & \textbf{Tác dụng} \\
\midrule
\texttt{separation} & khoảng cách giữa cặp tâm gần nhất, tính theo bán kính cụm.
Núm độ khó chính: $1$ chồng lấn nặng, $6$ tầm thường \\
\texttt{radius\_range}, \texttt{weight\_range} & cụm không đều về kích thước và mật độ \\
\texttt{noise}, \texttt{noise\_in\_cluster} & nhiễu nền, nhãn $\bot$ nên metric bỏ qua \\
\texttt{drift\_speed}, \texttt{drift\_every} & tâm cụm trôi và dội lại ở biên hypercube \\
\texttt{event\_frequency}, \texttt{events} & cụm gộp, tách, sinh ra, biến mất \\
\texttt{block\_ordered} & \textbf{control}: phát theo khối cùng lớp \\
\bottomrule
\end{tabular}
\end{table}

\begin{remark}[Vì sao \texttt{separation} phải \emph{suy ra} bán kính]
\label{rem:separation}
Cách hiển nhiên là đặt \texttt{separation} làm ràng buộc khoảng cách tối thiểu
khi rải tâm. Cách đó \emph{không hoạt động}: hypercube càng rộng khi $d$ tăng,
nên với $d = 10$ các tâm rải đều luôn cách nhau khoảng $1.2$ trong khi ràng buộc
chỉ đòi $0.1$--$0.4$; ràng buộc không bao giờ bị chạm và mọi layout đều dễ như
nhau. Đo được: \texttt{separation=0.8} và \texttt{separation=3.0} cho kết quả
\emph{giống hệt} từng chữ số. Framework vì thế suy ra bán kính từ layout,
$r = d_{\min} / (2\,\texttt{separation})$, khiến ràng buộc đúng theo dựng hình và
độ khó so sánh được giữa các $d$ và các $k$.
\end{remark}

\noindent
Bảy kịch bản đặt tên sẵn: \texttt{static\_easy}, \texttt{static\_overlap},
\texttt{unbalanced}, \texttt{noisy}, \texttt{incremental\_drift},
\texttt{evolving}, \texttt{block\_ordered}. Bảng~\ref{tab:generator} liệt kê các
núm điều khiển, và Nhận xét~\ref{rem:separation} giải thích vì sao núm độ khó
phải suy ra bán kính thay vì chặn khoảng cách.

\subsection{Luồng thực tế}

Cột $\rho_{\text{deg}}$ là tỉ lệ cửa sổ $1000$ mẫu chỉ chứa một lớp, đo trên
$50\,000$ bản ghi đầu. Nó không phải điểm chất lượng của dataset mà là
\emph{điều kiện đọc}: $\rho_{\text{deg}}$ cao nghĩa là chỉ số per-window thoái
hoá ở đó (Mục~\ref{sec:degenerate}) và chỉ nên đọc con số toàn stream.

\begin{table}[htbp]
\centering
\small
\caption{Các luồng thực tế framework nạp sẵn.}
\label{tab:realdata}
\begin{tabular}{@{}lrrrrr@{}}
\toprule
\textbf{Dataset} & $n$ & $d$ & \textbf{lớp} & $\rho_{\text{deg}}$ & \textbf{trung vị lớp áp đảo} \\
\midrule
\texttt{kdd99} & 494k & 34 & 23 & $\mathbf{72\%}$ & 1.000 \\
\texttt{covtype} & 581k & 10 & 7 & $0\%$ & 0.735 \\
\texttt{insects} & 24k--355k & 33 & 6 & $2\%$ & \textbf{0.276} \\
\texttt{keystroke} & 20k & 31 & 51 & $0\%$ & 0.400 \\
\texttt{shuttle} & 49k & 9 & 2 & $0\%$ & 0.931 \\
\texttt{electricity} & 45k & 8 & 2 & $0\%$ & 0.579 \\
\texttt{sms\_spam} & 5.5k & text & 2 & --- & --- \\
\bottomrule
\end{tabular}
\end{table}

\texttt{sms\_spam} có mặt vì không phải mọi clusterer của River đều ăn véc-tơ toạ
độ: TextClust cần token counts, nên cần một luồng văn bản và bước
\texttt{bag\_of\_words}.

\paragraph{Thêm dataset của bạn.}
\begin{lstlisting}[language=Python]
from benchmark.datasets import IterableDataset

def load_mine(path):
    def factory():                     # called fresh on every iteration
        for row in read(path):
            yield {"a": row.a, "b": row.b}, row.label
    return IterableDataset(factory, name="mine", n_features=2, has_labels=True)
\end{lstlisting}

\noindent
Với một dataset của River, \texttt{from\_river\_dataset()} lo phần còn lại --- kể
cả các dataset mà cột số đến dưới dạng \emph{chuỗi}, một chi tiết đủ để một bộ
lọc \texttt{isinstance(v, Number)} loại sạch mọi đặc trưng và trả về luồng rỗng.

% =====================================================================
\section{Giao thức đánh giá}
\label{sec:protocol}
% =====================================================================

\subsection{Điều kiện so sánh công bằng}

Mọi setting phải chạy trong điều kiện giống hệt nhau: cùng dataset, cùng thứ tự
luồng, cùng số mẫu, cùng lưới cửa sổ đánh giá, cùng phần cứng, cùng tập seed.
Framework cưỡng chế điều này bằng cách cho mọi setting dùng chung một đối tượng
dataset trả về iterator mới ở mỗi lần lặp, và ghi lại toàn bộ siêu tham số ---
\emph{kể cả tham số của scaler} --- vào kết quả.

Về siêu tham số, chúng tôi tuân theo nguyên tắc không tinh chỉnh trên chính tập
kiểm thử:
\[
\text{luồng hiệu chuẩn} \longrightarrow \text{tinh chỉnh} \longrightarrow
\text{đóng băng tham số} \longrightarrow \text{luồng kiểm thử} \longrightarrow
\text{benchmark cuối}.
\]
Trong mã nguồn, các tham số đã đóng băng nằm trong hằng \texttt{PRESETS} ở đầu
script benchmark.

\subsection{Hai họ chỉ số chất lượng: luỹ kế và theo cửa sổ}
\label{sec:cumwin}

Chia luồng thành các cửa sổ liên tiếp $W_1, W_2, \dots, W_K$ kích thước $m$
(mặc định $m=1000$), $W_k = \{(k-1)m+1,\dots,km\}$. Với một độ đo $M$ bất kỳ,
framework theo dõi \emph{đồng thời} hai đại lượng:

\begin{definition}[Chỉ số luỹ kế]
\begin{equation}
  M^{\mathrm{cum}}_k \;=\; M\Big(\{(y_t,\hat y_t)\}_{t \le km}\Big).
\end{equation}
Đối tượng metric không bao giờ được reset. Dãy này hội tụ, và
$M^{\mathrm{cum}}_K$ chính là điểm số toàn luồng.
\end{definition}

\begin{definition}[Chỉ số theo cửa sổ]
\begin{equation}
  M^{\mathrm{win}}_k \;=\; M\Big(\{(y_t,\hat y_t)\}_{t \in W_k}\Big).
\end{equation}
Đối tượng metric được khởi tạo lại ở mỗi biên cửa sổ, nên $M^{\mathrm{win}}_k$
chỉ phụ thuộc vào các mẫu \emph{bên trong} $W_k$.
\end{definition}

Mỗi họ được tóm tắt bằng ba con số:
\begin{equation}
  \label{eq:summary}
  \operatorname{mean}(M) = \frac{1}{K}\sum_{k=1}^{K} M_k,
  \qquad
  \operatorname{median}(M),
  \qquad
  \operatorname{final}(M) = M_K.
\end{equation}

\begin{proposition}[Trung bình đường luỹ kế không phải trung bình theo thời gian]
\label{prop:cumbias}
Nếu $M$ là trung bình cộng của một đại lượng theo mẫu, thì
\begin{equation}
  M^{\mathrm{cum}}_k \;=\; \frac{1}{k}\sum_{j=1}^{k} M^{\mathrm{win}}_j,
\end{equation}
và do đó
\begin{equation}
  \frac{1}{K}\sum_{k=1}^{K} M^{\mathrm{cum}}_k
  \;=\; \sum_{j=1}^{K} w_j\, M^{\mathrm{win}}_j,
  \qquad
  w_j \;=\; \frac{1}{K}\sum_{k=j}^{K}\frac{1}{k},
\end{equation}
với các trọng số $w_j$ \emph{giảm nghiêm ngặt} theo $j$.
\end{proposition}

\begin{proof}
Đổi thứ tự lấy tổng:
$\sum_{k=1}^{K}\frac{1}{k}\sum_{j=1}^{k}M^{\mathrm{win}}_j
= \sum_{j=1}^{K} M^{\mathrm{win}}_j \sum_{k=j}^{K}\frac1k$.
Vì $\sum_{k=j}^{K}1/k$ giảm khi $j$ tăng, ta có $w_1 > w_2 > \cdots > w_K$.
\end{proof}

Nói cách khác, lấy trung bình đường luỹ kế đặt trọng số cao nhất lên cửa sổ
\emph{đầu tiên} --- đúng giai đoạn khởi động lạnh, khi cả scaler lẫn mô hình đều
chưa ổn định. Đây không phải là ``chất lượng trung bình dọc luồng''; nó là một
ước lượng bị kéo lệch về phía đầu luồng. Trên luồng dừng, hệ quả là
$\operatorname{mean}(M^{\mathrm{cum}})$ \emph{đánh giá thấp} chất lượng ổn định
của mô hình --- điều chúng tôi đo được ở Bảng~\ref{tab:cumwin}. Vì vậy chỉ số
time-aware mà framework báo cáo mặc định là $\operatorname{mean}(M^{\mathrm{win}})$.
Đường luỹ kế vẫn được giữ và báo cáo riêng, vì nó là đường hội tụ và
$\operatorname{final}(M^{\mathrm{cum}})$ chính là điểm toàn luồng.

\subsection{Độ đo ngoại (cần ground truth)}
\label{sec:external}

Cho bảng liên hợp (contingency table) giữa $R$ lớp thật và $C$ cụm dự đoán trên
tập mẫu kích thước $n$:
\begin{equation}
  n_{ij} = \big|\{t : y_t = i \wedge \hat y_t = j\}\big|,
  \qquad
  a_i = \sum_j n_{ij},
  \qquad
  b_j = \sum_i n_{ij}.
\end{equation}

\paragraph{Adjusted Rand Index.}
River tính ARI qua ma trận nhầm lẫn theo \emph{cặp} điểm. Đặt
\begin{align}
  \mathrm{TP} &= \textstyle\sum_{ij} n_{ij}^2 - n
    &&\text{(cặp cùng lớp, cùng cụm)},\\
  \mathrm{FP} &= \textstyle\sum_{j} b_j^2 - \sum_{ij} n_{ij}^2
    &&\text{(cặp khác lớp, cùng cụm)},\\
  \mathrm{FN} &= \textstyle\sum_{i} a_i^2 - \sum_{ij} n_{ij}^2
    &&\text{(cặp cùng lớp, khác cụm)},\\
  \mathrm{TN} &= n^2 - (\mathrm{FP}+\mathrm{FN}) - \textstyle\sum_{ij} n_{ij}^2
    &&\text{(cặp khác lớp, khác cụm)},
\end{align}
(tất cả đếm theo cặp có thứ tự). Khi đó
\begin{equation}
  \label{eq:ari}
  \boxed{\;
  \mathrm{ARI} \;=\;
  \frac{2\big(\mathrm{TP}\cdot\mathrm{TN} \;-\; \mathrm{FN}\cdot\mathrm{FP}\big)}
       {(\mathrm{TP}+\mathrm{FN})(\mathrm{FN}+\mathrm{TN})
        \;+\;(\mathrm{TP}+\mathrm{FP})(\mathrm{FP}+\mathrm{TN})}
  \;}
\end{equation}
với quy ước: nếu mẫu số bằng $0$ thì $\mathrm{ARI} = 1$. Giá trị $1$ nghĩa là
trùng khớp hoàn hảo, $0$ nghĩa là bằng mức ngẫu nhiên, và ARI \emph{có thể âm}
khi phân hoạch còn tệ hơn ngẫu nhiên.

\paragraph{Mutual Information và các biến thể chuẩn hoá.}
\begin{align}
  I(U,V) &= \sum_{i=1}^{R}\sum_{j=1}^{C}
    \frac{n_{ij}}{n}\,\log\!\frac{n\,n_{ij}}{a_i\,b_j},
    \label{eq:mi}\\
  H(U) &= -\sum_{i=1}^{R} \frac{a_i}{n}\log\frac{a_i}{n},
  \qquad
  H(V) = -\sum_{j=1}^{C} \frac{b_j}{n}\log\frac{b_j}{n},\\
  \mathrm{NMI} &= \frac{I(U,V)}{\tfrac{1}{2}\big(H(U)+H(V)\big)},
    \label{eq:nmi}\\
  \mathrm{AMI} &= \frac{I(U,V) - \E[I(U,V)]}
                       {\tfrac{1}{2}\big(H(U)+H(V)\big) - \E[I(U,V)]},
    \label{eq:ami}
\end{align}
trong đó chuẩn hoá mặc định là trung bình cộng của hai entropy, và $\E[I]$ là kỳ
vọng của thông tin tương hỗ dưới mô hình siêu bội với các biên $\{a_i\},\{b_j\}$
cố định. Công thức~\eqref{eq:mi} là thông tin tương hỗ,~\eqref{eq:nmi} và~\eqref{eq:ami}
là hai cách chuẩn hoá nó. NMI \emph{không} hiệu chỉnh theo may rủi nên tăng một
cách máy móc theo số cụm; AMI thì có. Các trường hợp biên được River quy ước như sau:
$R=C=1 \Rightarrow \mathrm{NMI}=\mathrm{AMI}=1$, và
$(R=1) \veebar (C=1) \Rightarrow \mathrm{AMI}=0$.

\paragraph{Purity.}
\begin{equation}
  \label{eq:purity}
  \mathrm{Purity} \;=\; \frac{1}{n}\sum_{j=1}^{C} \max_{i} n_{ij}.
\end{equation}
Purity nằm trong $[0,1]$ nhưng \emph{không} hiệu chỉnh theo may rủi: nó tăng đơn
điệu theo số cụm và đạt $1$ khi mỗi mẫu là một cụm riêng. Nó cũng bị thổi phồng
khi phân bố lớp mất cân bằng nặng --- trên KDDCup99, gán mọi thứ vào một cụm duy
nhất vẫn cho purity $\approx 0.80$. Vì vậy purity chỉ nên đọc kèm ARI/AMI.

\subsection{Độ đo nội (không cần ground truth)}

River cài đặt một phiên bản streaming của hệ số Silhouette dựa trên tâm cụm chứ
không phải khoảng cách từng cặp,~\eqref{eq:silhouette}:
\begin{equation}
  \label{eq:silhouette}
  \mathrm{Sil}^{\text{stream}}
  \;=\;
  \frac{\sum_{t} \big\| \tilde\xt_t - \mathbf{c}_{\hat y_t} \big\|_2}
       {\sum_{t} \big\| \tilde\xt_t - \mathbf{c}^{(2)}(\tilde\xt_t) \big\|_2},
  \qquad
  \mathbf{c}^{(2)}(\xt) = \operatorname*{arg\,min}_{\mathbf{c} \neq \mathbf{c}_{\hat y}} \|\xt - \mathbf{c}\|_2 .
\end{equation}

\begin{remark}[Hai cái bẫy của Silhouette streaming]
Thứ nhất, đây là tỉ số \emph{cố kết / phân tách}, nên \textbf{nhỏ hơn là tốt
hơn} --- ngược chiều với silhouette theo lô của scikit-learn.
Thứ hai, tử số và mẫu số phải được tính trong \emph{cùng không gian} mà mô hình
làm việc: tâm $\mathbf{c}$ sống trong không gian đã chuẩn hoá, nên phải dùng
$\tilde\xt_t$ chứ không phải $\xt_t$. Đưa nhầm điểm thô vào làm chỉ số sai lệch
hơn một bậc độ lớn (chúng tôi quan sát $0.90$ so với $0.04$ trên cùng một lần
chạy).
\end{remark}

\subsection{Chỉ số hiệu năng}
\label{sec:efficiency}

Gọi $\tau_t$ là thời gian tường (wall time) của bước~\eqref{eq:ttt} tại mẫu $t$,
đo bằng đồng hồ đơn điệu bao quanh \emph{chỉ} phần scale--predict--learn. Chi
phí đọc dữ liệu và cập nhật metric nằm ngoài phép đo, vì chúng không thuộc về
thuật toán.
\begin{equation}
  T_{\text{total}} = \sum_{t=1}^{n}\tau_t,
  \qquad
  \text{throughput} = \frac{n}{T_{\text{total}}}\ \text{[mẫu/s]},
  \qquad
  \text{ms/mẫu} = \frac{1000\,T_{\text{total}}}{n}.
\end{equation}

Về bộ nhớ, framework báo hai đại lượng khác nhau về bản chất:
\begin{itemize}
  \item \textbf{Bộ nhớ mô hình} --- kích thước đồ thị đối tượng của bộ phân cụm
        cộng scaler, lấy qua \texttt{river.base.Base.\_raw\_memory\_usage}. Đây
        là con số \emph{so sánh được} giữa các thuật toán: nó không tính
        interpreter, dataset hay các đối tượng metric.
  \item \textbf{RSS tiến trình} --- bộ nhớ thường trú của cả tiến trình. Chỉ để
        tham chiếu như một chặn trên.
\end{itemize}

\begin{remark}[Vì sao chỉ đo bộ nhớ tại checkpoint]
\label{rem:memcost}
\texttt{\_raw\_memory\_usage} duyệt toàn bộ đồ thị đối tượng bằng
\texttt{gc.get\_referents}, nên chi phí của nó là $\mathcal{O}(|\theta|)$ --- tỉ
lệ với \emph{kích thước mô hình}, không phải hằng số. Với một bộ phân cụm giữ
đệm $2000$ điểm, một lần gọi tốn khoảng $14$\,ms, tức gấp hơn hai bậc độ lớn so
với chi phí xử lý một mẫu. Gọi nó mỗi mẫu sẽ biến một benchmark $400\,000$ mẫu
thành hàng giờ đồng hồ thuần tuý cho việc \emph{đo}. Vì vậy phép đo bộ nhớ chỉ
chạy tại các checkpoint, và hệ quả là
\begin{equation}
  \text{peak memory} = \max_{\text{checkpoint } c} |\theta_c|
\end{equation}
là một \emph{chặn dưới} của đỉnh thật. Với mô hình có tỉa (DenStream, DBSTREAM),
cần tăng số checkpoint nếu muốn chặn chặt hơn.
\end{remark}

Số micro-cluster $|\mathcal{M}_t|$ cũng được lấy mẫu tại checkpoint; nó là biến
giải thích chính cho cả runtime lẫn bộ nhớ, vì hầu hết các thuật toán có chi phí
mỗi mẫu tăng theo $|\mathcal{M}_t|$ (DBSTREAM thậm chí có bước dọn dẹp
$\mathcal{O}(|\mathcal{M}|^2)$).

\subsection{Độ ổn định qua nhiều seed}

Với thuật toán có thành phần ngẫu nhiên, mỗi cấu hình được chạy lại trên tập
seed $\mathcal{Z}$ và báo cáo dưới dạng
\begin{equation}
  \bar{M} \pm s_M,
  \qquad
  \bar{M} = \frac{1}{|\mathcal{Z}|}\sum_{z\in\mathcal{Z}} M^{(z)},
  \qquad
  s_M = \sqrt{\frac{1}{|\mathcal{Z}|}\sum_{z\in\mathcal{Z}}\big(M^{(z)}-\bar M\big)^2}.
\end{equation}

Framework xác định tính ngẫu nhiên bằng cách nội quan chữ ký hàm khởi tạo: nếu
lớp có tham số \texttt{seed} (hoặc có \texttt{**kwargs} chuyển tiếp xuống một mô
hình con có \texttt{seed}) thì seed được tiêm vào. Thuật toán \emph{tất định}
(DenStream, DBSTREAM) chỉ chạy một lần, vì lặp lại chúng trên cùng một luồng cho
đúng cùng con số; báo cáo ghi rõ số lần chạy thật của từng setting để độ lệch
chuẩn bằng $0$ không bị đọc nhầm là ``rất ổn định''.

% =====================================================================
\section{Cửa sổ thoái hoá: khi chỉ số time-aware mất ý nghĩa}
\label{sec:degenerate}
% =====================================================================

Đây là phần đóng góp phương pháp chính của báo cáo.

\begin{definition}[Cửa sổ thoái hoá]
Cửa sổ $W_k$ được gọi là \emph{thoái hoá} nếu tập nhãn thật bên trong nó chỉ có
một phần tử: $\big|\{y_t : t \in W_k\}\big| = 1$.
\end{definition}

\begin{proposition}[ARI suy biến trên cửa sổ thoái hoá]
\label{prop:degenerate}
Trên một cửa sổ thoái hoá, giá trị ARI theo~\eqref{eq:ari} chỉ nhận hai giá trị:
\begin{equation}
  \mathrm{ARI}\big(W_k\big) =
  \begin{cases}
    1, & \text{nếu mô hình gán toàn bộ cửa sổ vào đúng một cụm},\\
    0, & \text{với mọi cách chia khác, bất kể chia thế nào}.
  \end{cases}
\end{equation}
\end{proposition}

\begin{proof}
Vì chỉ có một lớp thật, mọi cặp điểm trong cửa sổ đều là cặp ``cùng lớp''. Do đó
không tồn tại cặp khác lớp, kéo theo
\[
  \mathrm{FP} = 0, \qquad \mathrm{TN} = 0 .
\]
Thay vào~\eqref{eq:ari}, tử số trở thành
$2(\mathrm{TP}\cdot 0 - \mathrm{FN}\cdot 0) = 0$, còn mẫu số là
\[
  (\mathrm{TP}+\mathrm{FN})\,\mathrm{FN} \;+\; (\mathrm{TP}+0)\cdot 0
  \;=\; (\mathrm{TP}+\mathrm{FN})\,\mathrm{FN}.
\]
Nếu mô hình chia cửa sổ thành từ hai cụm trở lên thì tồn tại cặp cùng lớp khác
cụm, tức $\mathrm{FN}>0$, nên mẫu số dương và $\mathrm{ARI}=0$. Nếu mô hình gộp
toàn bộ vào một cụm thì $\mathrm{FN}=0$, mẫu số bằng $0$, và theo quy ước
$0/0 \mapsto 1$ ta được $\mathrm{ARI}=1$.
\end{proof}

\begin{remark}
Kết quả tương tự đúng cho AMI theo quy ước biên của River: $R=1$ cho
$\mathrm{AMI}=0$ trừ khi $C=1$, lúc đó $\mathrm{AMI}=1$. Purity thì luôn bằng
$1$ trên cửa sổ thoái hoá, bất kể mô hình làm gì.
\end{remark}

Hệ quả thực tiễn rất nghiêm trọng: trên cửa sổ thoái hoá, chỉ số per-window
không đo chất lượng phân cụm mà đo \emph{mô hình có gộp cửa sổ lại hay không}.
Một mô hình \emph{phân mảnh quá ít} (under-segmenting) --- thường là mô hình
tồi --- sẽ ăn trọn điểm $1.0$ trên mọi cửa sổ thoái hoá. Nếu luồng đến theo từng
khối cùng lớp, trung bình per-window vì thế trở thành một phần thưởng cho việc
phân cụm kém.

\paragraph{Thống kê chẩn đoán.}
Framework do đó tính và báo cáo, cho mỗi lần chạy:
\begin{align}
  \rho_{\text{deg}} &= \frac{1}{K}\sum_{k=1}^{K}
    \ind{\,\big|\{y_t : t \in W_k\}\big| = 1\,}
    &&\text{(tỉ lệ cửa sổ thoái hoá)},\\
  \tilde{p} &= \operatorname{median}_{k}
    \left( \frac{\max_i |\{t\in W_k : y_t = i\}|}{|W_k|} \right)
    &&\text{(trung vị tỉ lệ lớp áp đảo)}.
\end{align}
Bảng chẩn đoán được in \emph{trước} bảng time-aware, kèm cảnh báo tự động khi
$\rho_{\text{deg}} \ge 0.3$.

% =====================================================================
\section{Mã giả}
\label{sec:pseudocode}
% =====================================================================

\begin{algorithm}[htbp]
\caption{\textsc{RunOnce} --- benchmark một setting trên một luồng}
\label{alg:runonce}
\begin{algorithmic}[1]
\Require luồng $\mathcal{S}$; adapter $A$ (scaler + bộ phân cụm); cỡ cửa sổ $m$;
         số checkpoint $P$; tập độ đo ngoại $\mathcal{E}$, nội $\mathcal{I}$
\Ensure  quỹ đạo chất lượng, quỹ đạo tài nguyên, chẩn đoán
\State $M^{\mathrm{cum}} \gets \textsc{NewMetrics}(\mathcal{E}\cup\mathcal{I})$
       \Comment{không bao giờ reset}
\State $M^{\mathrm{win}} \gets \textsc{NewMetrics}(\mathcal{E}\cup\mathcal{I})$
       \Comment{reset mỗi cửa sổ}
\State $L \gets \emptyset$ \Comment{đa tập nhãn của cửa sổ hiện tại}
\State $g \gets \lfloor |\mathcal{S}| / P \rfloor$;\quad
       $T \gets 0$;\quad $n \gets 0$
\ForAll{$(\xt, y) \in \mathcal{S}$}
  \If{$n > 0$ \textbf{and} $n \bmod m = 0$}
      \Comment{biên cửa sổ}
    \State \textsc{FlushWindow}$(M^{\mathrm{cum}}, M^{\mathrm{win}}, L)$
    \State $M^{\mathrm{win}} \gets \textsc{NewMetrics}(\mathcal{E}\cup\mathcal{I})$;\quad
           $L \gets \emptyset$
  \EndIf
  \State $t_0 \gets \textsc{Clock}()$
  \State $\tilde\xt \gets A.\textsc{Scale}(\xt)$
         \Comment{learn rồi mới transform, \eqref{eq:scaleorder}}
  \State $\hat y \gets A.\textsc{Predict}(\tilde\xt)$
         \Comment{test truớc}
  \State $A.\textsc{Learn}(\tilde\xt)$
         \Comment{train sau}
  \State $T \gets T + (\textsc{Clock}() - t_0)$;\quad $n \gets n+1$
  \If{$y \neq \bot$}
    \State $L \gets L \uplus \{y\}$
    \State cập nhật mọi độ đo ngoại trong $M^{\mathrm{cum}}$ và $M^{\mathrm{win}}$ bằng $(y,\hat y)$
  \EndIf
  \If{$\mathcal{I} \neq \emptyset$ \textbf{and} $A.\textsc{Centers}() \neq \emptyset$}
    \State cập nhật mọi độ đo nội bằng $(\tilde\xt, \hat y, A.\textsc{Centers}())$
           \Comment{cùng không gian với tâm cụm}
  \EndIf
  \If{$n \bmod g = 0$}
    \State \textsc{Checkpoint}$(n, T)$ \Comment{Thuật toán~\ref{alg:checkpoint}}
  \EndIf
\EndFor
\State \textsc{FlushWindow}$(\cdot)$ cho cửa sổ dở dang cuối; \textsc{Checkpoint}$(n,T)$
\State \Return tóm tắt~\eqref{eq:summary} trên $\{M^{\mathrm{win}}_k\}$ \emph{và} trên
       $\{M^{\mathrm{cum}}_k\}$, cùng $\rho_{\text{deg}}, \tilde p$
\end{algorithmic}
\end{algorithm}

\begin{algorithm}[htbp]
\caption{\textsc{Checkpoint} --- lấy mẫu tài nguyên (đắt, nên thưa)}
\label{alg:checkpoint}
\begin{algorithmic}[1]
\Require số mẫu đã xử lý $n$, runtime luỹ kế $T$
\State $\mu \gets A.\textsc{MicroClusterCount}()$;\quad
       $\kappa \gets A.\textsc{MacroClusterCount}()$
\State $b \gets A.\textsc{ModelMemoryBytes}()$
       \Comment{$\mathcal{O}(|\theta|)$ --- xem Nhận xét~\ref{rem:memcost}}
\State $\varrho \gets \textsc{ProcessRSS}()$
\State $\text{peak}_{\text{mem}} \gets \max(\text{peak}_{\text{mem}}, b)$;\quad
       $\text{peak}_{\text{rss}} \gets \max(\text{peak}_{\text{rss}}, \varrho)$
\State ghi bộ $(n,\ T,\ b,\ \varrho,\ \mu,\ \kappa)$
\end{algorithmic}
\end{algorithm}

\begin{algorithm}[htbp]
\caption{\textsc{RunBenchmark} --- nhiều setting, nhiều seed, cùng một luồng}
\label{alg:harness}
\begin{algorithmic}[1]
\Require danh sách setting $\mathcal{P}$; dataset $D$; tập seed $\mathcal{Z}$
\State ghi lại môi trường: CPU, RAM, OS, phiên bản Python và River
\ForAll{$p \in \mathcal{P}$}
  \If{\textsc{AcceptsSeed}$(p.\text{algorithm})$}
     \Comment{nội quan chữ ký hàm khởi tạo}
    \State $Z_p \gets \mathcal{Z}$
  \Else
    \State $Z_p \gets \{\,\text{seed đầu tiên}\,\}$
           \Comment{tất định: lặp lại là lãng phí}
  \EndIf
  \ForAll{$z \in Z_p$}
    \State $A \gets$ \textsc{BuildAdapter}$(p, z)$
           \Comment{tham số scaler cũng được ghi lại}
    \State $r_{p,z} \gets$ \textsc{RunOnce}$(D, A, \dots)$
           \Comment{$D$ trả iterator mới $\Rightarrow$ cùng luồng}
  \EndFor
  \State $\text{agg}_p \gets$ mean $\pm$ std trên $\{r_{p,z}\}_{z \in Z_p}$
\EndFor
\State \Return bảng, đồ thị, JSON, kèm số lần chạy thật $|Z_p|$ của từng setting
\end{algorithmic}
\end{algorithm}

Thuật toán~\ref{alg:runonce} là vòng lặp trực tuyến chính,
Thuật toán~\ref{alg:checkpoint} là bước lấy mẫu tài nguyên, và
Thuật toán~\ref{alg:harness} là lớp điều phối nhiều setting--nhiều seed bên trên.

Thuật toán~\ref{alg:runonce} là vòng lặp trực tuyến chính,
Thuật toán~\ref{alg:checkpoint} là bước lấy mẫu tài nguyên đắt tiền, và
Thuật toán~\ref{alg:harness} là lớp điều phối nhiều setting, nhiều seed bên
trên. Nhận xét~\ref{rem:guard} giải thích vì sao bước chuẩn hoá ở dòng~11 của
Thuật toán~\ref{alg:runonce} cần một điều kiện bảo vệ chặt hơn vẻ ngoài của nó.

% =====================================================================
\section{Kiểm chứng framework}
\label{sec:results}
% =====================================================================

Các số đo dưới đây tồn tại để \emph{kiểm chứng bộ máy đo}, không phải để xếp
hạng thuật toán hay đề xuất siêu tham số. Chúng trả lời những câu hỏi về chính
phép đo: chi phí của một phép đo có làm hỏng thứ nó đo không
(Mục~\ref{sec:memcost-res}); hai họ chỉ số chất lượng có thực sự khác nhau không
(Mục~\ref{sec:cumwin-res}); Mệnh đề~\ref{prop:degenerate} có đúng trên dữ liệu
thật không (Mục~\ref{sec:degen-res}). Mọi cấu hình thuật toán xuất hiện ở đây là
\emph{ví dụ chạy được}, không phải khuyến nghị. Các luồng dữ liệu dùng ở đây mô
tả ở Mục~\ref{sec:datasets} (Bảng~\ref{tab:generator} và~\ref{tab:realdata}).

\subsection{Môi trường}

Bảng~\ref{tab:env} là bản ghi môi trường áp dụng cho mọi số đo trong mục này;
framework tự sinh nó ở mỗi lần chạy.

\begin{table}[htbp]
\centering
\small
\caption{Môi trường chạy, do framework tự ghi lại ở mỗi lần benchmark.}
\label{tab:env}
\begin{tabular}{@{}ll@{}}
\toprule
CPU & AMD64 Family 25 Model 80 (8 nhân vật lý / 16 luồng) \\
RAM & 27.86 GB \\
OS & Windows 11 (10.0.26200) \\
Python & 3.13.3 (CPython) \\
River & 0.26.1 (cây nguồn cục bộ); bản phân phối đã cài: 0.26.0 \\
NumPy / SciPy / scikit-learn & 2.4.3 / 1.17.1 / 1.8.0 \\
\bottomrule
\end{tabular}
\end{table}

Chênh lệch giữa \texttt{0.26.1} (phiên bản trong cây nguồn đang được import) và
\texttt{0.26.0} (siêu dữ liệu gói đã cài) chính là loại sai lệch mà bản ghi môi
trường sinh ra để phát hiện.

\subsection{Chi phí của phép đo bộ nhớ}
\label{sec:memcost-res}

Bảng~\ref{tab:memcost} đo trực tiếp chi phí nêu ở Nhận xét~\ref{rem:memcost}:
thời gian xử lý $3000$ mẫu $34$ chiều, so với $3000$ lần đọc
\texttt{\_raw\_memory\_usage}.

\begin{table}[htbp]
\centering
\small
\caption{Chi phí đo bộ nhớ so với chi phí xử lý, trên $3000$ mẫu.}
\label{tab:memcost}
\begin{tabular}{@{}lrrr@{}}
\toprule
\textbf{Thuật toán} & \textbf{Xử lý (s)} & \textbf{$3000$ lần đo (s)} & \textbf{Hệ số} \\
\midrule
DBSTREAM      & 0.104 & 0.207  & $2.0\times$ \\
STREAMKMeans  & 0.241 & 42.771 & $\mathbf{177.6\times}$ \\
\bottomrule
\end{tabular}
\end{table}

Ngoại suy sang một benchmark $400\,000$ mẫu, riêng phép đo bộ nhớ của
STREAMKMeans sẽ tốn khoảng $1.6$ giờ --- nhiều hơn toàn bộ phần tính toán thật
hàng trăm lần. Đây là lý do phép đo được dời về checkpoint.

\subsection{Luỹ kế so với theo cửa sổ}
\label{sec:cumwin-res}

\begin{table}[htbp]
\centering
\small
\caption{Hai cách tóm tắt ARI trên luồng Gaussian tổng hợp
($n=20\,000$, $d=10$, $k=5$, $m=1000$). Trung bình đường luỹ kế \emph{đánh giá
thấp} chất lượng ổn định, đúng như Mệnh đề~\ref{prop:cumbias} dự đoán: nó mang
theo giai đoạn khởi động lạnh với trọng số lớn nhất.}
\label{tab:cumwin}
\begin{tabular}{@{}llrrrr@{}}
\toprule
\textbf{Thuật toán} & \textbf{Seed}
 & $\operatorname{mean}(M^{\mathrm{cum}})$ & $\operatorname{final}(M^{\mathrm{cum}})$
 & $\operatorname{mean}(M^{\mathrm{win}})$ & $\operatorname{final}(M^{\mathrm{win}})$ \\
\midrule
CluStream     & 1 & 0.868 & 0.958 & \textbf{0.965} & 1.000 \\
CluStream     & 2 & 0.177 & 0.175 & \textbf{0.299} & 0.278 \\
DenStream     & 1 & 0.092 & 0.084 & \textbf{0.198} & 0.201 \\
DBSTREAM      & 1 & 0.959 & 0.988 & \textbf{0.989} & 1.000 \\
STREAMKMeans  & 1 & 0.657 & 0.746 & \textbf{0.750} & 0.777 \\
STREAMKMeans  & 2 & 0.386 & 0.444 & \textbf{0.457} & 0.453 \\
\bottomrule
\end{tabular}
\end{table}

Trong cả sáu lần chạy, $\operatorname{mean}(M^{\mathrm{cum}}) <
\operatorname{mean}(M^{\mathrm{win}})$, độ lệch từ $0.03$ đến $0.12$ điểm ARI.
Hình~\ref{fig:cumwin} vẽ hai đường cho CluStream (seed 1): đường per-window đạt
$1.0$ sau cửa sổ thứ ba và giữ nguyên, còn đường luỹ kế mất tới cuối luồng mới
bò lên $0.958$ vì nó vẫn đang gánh $500$ mẫu đầu tiên.

\begin{figure}[htbp]
\centering
\begin{tikzpicture}
\begin{axis}[
  width=13cm, height=6.2cm,
  xlabel={số mẫu đã xử lý $t$},
  ylabel={ARI},
  ymin=0.3, ymax=1.06,
  xmin=0, xmax=20000,
  legend pos=south east,
  legend cell align=left,
  grid=both,
  grid style={gray!18},
  tick label style={font=\small},
  label style={font=\small},
  legend style={font=\small},
]
\addplot[thick,blue!70!black,mark=*,mark size=1.2pt] coordinates {
(499,0.4315)(1499,0.6305)(2499,0.7403)(3499,0.8012)(4499,0.8380)(5499,0.8640)
(6499,0.8827)(7499,0.8971)(8499,0.9083)(9499,0.9170)(10499,0.9246)(11499,0.9308)
(12499,0.9360)(13499,0.9404)(14499,0.9444)(15499,0.9478)(16499,0.9508)
(17499,0.9535)(18499,0.9560)(19499,0.9582)
};
\addlegendentry{luỹ kế (cumulative)}
\addplot[thick,orange!85!black,mark=square*,mark size=1.2pt] coordinates {
(499,0.4315)(1499,0.8644)(2499,1.0000)(3499,1.0000)(4499,1.0000)(5499,1.0000)
(6499,1.0000)(7499,1.0000)(8499,1.0000)(9499,1.0000)(10499,1.0000)(11499,1.0000)
(12499,1.0000)(13499,1.0000)(14499,1.0000)(15499,1.0000)(16499,1.0000)
(17499,1.0000)(18499,1.0000)(19499,1.0000)
};
\addlegendentry{theo cửa sổ (per-window)}
\end{axis}
\end{tikzpicture}
\caption{CluStream trên luồng tổng hợp, seed 1. Đường luỹ kế là đường
\emph{hội tụ}; đường per-window là đường \emph{tức thời}. Tóm tắt đường luỹ kế
rồi gọi đó là ``chất lượng trung bình dọc luồng'' sẽ báo $0.868$ trong khi mô
hình thực tế đã hoàn hảo từ mẫu thứ $2000$.}
\label{fig:cumwin}
\end{figure}

\subsection{Kiểm chứng thực nghiệm
  \texorpdfstring{Mệnh đề~\ref{prop:degenerate}}{Mệnh đề 2}}
\label{sec:degen-res}

Hình~\ref{fig:degenerate} là bằng chứng trực tiếp. Nó vẽ ARI per-window của
DBSTREAM trên $20\,000$ mẫu đầu KDDCup99 với hai scaler khác nhau, và tô nền
những cửa sổ thoái hoá.

\begin{figure}[htbp]
\centering
\begin{tikzpicture}
\begin{axis}[
  width=14cm, height=6.4cm,
  xlabel={số mẫu đã xử lý $t$},
  ylabel={ARI trong cửa sổ},
  ymin=-0.08, ymax=1.12,
  xmin=0, xmax=20000,
  legend pos=outer north east,
  legend cell align=left,
  grid=both, grid style={gray!18},
  tick label style={font=\small},
  label style={font=\small},
  legend style={font=\scriptsize},
]
% shade degenerate windows (n_true_classes == 1)
\pgfplotsinvokeforeach{1000,2000,3000,5000,6000,8000,9000,10000,12000,13000,14000,16000,17000,18000}{
  \addplot[draw=none,fill=red!9,forget plot] coordinates
    {(#1,-0.08) (\the\numexpr#1+1000\relax,-0.08)
     (\the\numexpr#1+1000\relax,1.12) (#1,1.12)} \closedcycle;
}
\addplot[thick,blue!70!black,mark=*,mark size=1.3pt] coordinates {
(499,0.0001)(1499,0.0000)(2499,0.0000)(3499,0.0000)(4499,0.0007)(5499,0.0000)
(6499,0.0000)(7499,0.1059)(8499,0.0000)(9499,0.0000)(10499,0.0000)(11499,0.4111)
(12499,0.0000)(13499,0.0000)(14499,0.0000)(15499,0.0150)(16499,0.0000)
(17499,0.0000)(18499,0.0000)(19499,0.0537)
};
\addlegendentry{\texttt{standard}}
\addplot[thick,orange!85!black,mark=square*,mark size=1.3pt] coordinates {
(499,0.0084)(1499,0.0000)(2499,1.0000)(3499,0.0000)(4499,0.0783)(5499,1.0000)
(6499,0.0000)(7499,0.0983)(8499,0.0000)(9499,1.0000)(10499,1.0000)(11499,0.0021)
(12499,0.0000)(13499,1.0000)(14499,1.0000)(15499,0.5080)(16499,0.0000)
(17499,1.0000)(18499,0.0000)(19499,0.8462)
};
\addlegendentry{\texttt{minmax}}
\end{axis}
\end{tikzpicture}
\caption{DBSTREAM trên KDDCup99 ($n=20\,000$, $m=1000$). Vùng nền hồng là các
cửa sổ \emph{thoái hoá} (chỉ chứa một lớp thật). Trên mọi cửa sổ đó, cả hai
đường chỉ nhận giá trị $0$ hoặc $1$ --- không có giá trị trung gian nào ---
đúng như Mệnh đề~\ref{prop:degenerate}. \texttt{minmax} nén dữ liệu mạnh nên gộp
được nhiều cửa sổ và ăn bảy điểm $1.0$ miễn phí; \texttt{standard} trải dữ liệu
ra nên luôn chia nhỏ và nhận $0$.}
\label{fig:degenerate}
\end{figure}

Trên $14$ cửa sổ thoái hoá, \texttt{minmax} nhận đúng bảy giá trị $1.0000$ và
bảy giá trị $0.0000$, không có giá trị nào khác; \texttt{standard} nhận $0.0000$
trên cả $14$. Không có ngoại lệ nào --- đúng như dự đoán lý thuyết.

Với KDDCup99 ở quy mô lớn hơn, mức độ thoái hoá còn cao hơn:

\begin{table}[htbp]
\centering
\small
\caption{Chẩn đoán cửa sổ trên KDDCup99 ($m=1000$).}
\label{tab:diag}
\begin{tabular}{@{}lrrrr@{}}
\toprule
$n$ & \# cửa sổ & \# thoái hoá & $\rho_{\text{deg}}$ & $\tilde p$ \\
\midrule
$20\,000$ & 20 & 14 & $70\%$ & 1.000 \\
$50\,000$ & 50 & 36 & $\mathbf{72\%}$ & 1.000 \\
\bottomrule
\end{tabular}
\end{table}

Bảng~\ref{tab:diag} tổng hợp ở hai quy mô.
Ở $n=50\,000$, ngoài $36$ cửa sổ thoái hoá còn có $40/50$ cửa sổ mà một lớp
chiếm $\ge 99\%$. Trung vị tỉ lệ lớp áp đảo bằng $1.000$.

\subsection{Khảo sát scaler: một cái bẫy hoàn chỉnh}
\label{sec:ablation}

Bảng~\ref{tab:ablation} khảo sát sáu biến thể chuẩn hoá với cùng một thuật toán
(DBSTREAM), cùng tham số, cùng luồng.

\begin{table}[htbp]
\centering
\small
\caption{Khảo sát scaler --- DBSTREAM trên KDDCup99, $n=50\,000$, seed 1.
Cột $\mathrm{ARI}^{\text{toàn}}$ là chỉ số toàn luồng; cột
$\overline{\mathrm{ARI}}^{\text{win}}$ là trung bình per-window.}
\label{tab:ablation}
\begin{tabular}{@{}lrrrrr@{}}
\toprule
\textbf{Scaler} & $\mathrm{ARI}^{\text{toàn}}$ & $\overline{\mathrm{ARI}}^{\text{win}}$
 & \textbf{mẫu/s} & \textbf{Model MB} & $\overline{|\mathcal{M}|}$ \\
\midrule
\texttt{standard}                       & $\mathbf{0.1458}$ & 0.0228 & 1\,022 & 0.66 & 109.5 \\
\texttt{standard(window\_size=10000)}   & 0.1225 & 0.0182 & 914 & $\mathbf{122.8}$ & 113.3 \\
\texttt{none}                           & 0.0491 & 0.0216 & $\mathbf{39\,860}$ & 0.05 & 4.9 \\
\texttt{robust}                         & 0.0039 & 0.0054 & 1\,692 & 0.81 & 140.0 \\
\texttt{adaptive\_standard(0.01)}       & $-0.0498$ & 0.0608 & 513 & 1.14 & 198.6 \\
\texttt{minmax}                         & $\mathbf{-0.0850}$ & $\mathbf{0.5377}$ & 4\,313 & 0.10 & 20.5 \\
\bottomrule
\end{tabular}
\end{table}

Bảng này chứa ba bài học tách biệt.

\paragraph{(a) Hai cột chất lượng cho hai bảng xếp hạng ngược nhau.}
Theo $\overline{\mathrm{ARI}}^{\text{win}}$, \texttt{minmax} thắng áp đảo
($0.5377$ so với $0.0228$, tức gấp $24$ lần). Theo $\mathrm{ARI}^{\text{toàn}}$,
\texttt{minmax} là biến thể \emph{tệ nhất} và điểm của nó \emph{âm} --- tệ hơn
gán ngẫu nhiên. Mệnh đề~\ref{prop:degenerate} cùng Hình~\ref{fig:degenerate}
giải thích trọn vẹn nghịch lý: điểm per-window của \texttt{minmax} gần như hoàn
toàn cấu thành từ những giá trị $1.0$ miễn phí trên cửa sổ thoái hoá. Trên luồng
kiểu này, \textbf{chỉ số toàn luồng mới là chỉ số nên tin}, và kết luận đúng là
dùng \texttt{standard}.

\paragraph{(b) Cái giá bộ nhớ của cửa sổ trượt.}
\texttt{standard(window\_size=10000)} tốn $122.8$ MB so với $0.66$ MB của
\texttt{standard} thường --- gấp $187$ lần. Điều này khớp với phân tích
$\mathcal{O}(dW)$: với $d=34$ và $W=10^4$, bộ chuẩn hoá phải lưu $3.4\times10^5$
quan sát. Đây là đánh đổi có thật, và nó chỉ hiện ra khi hàm tạo scaler cho phép
truyền tham số.

\paragraph{(c) Chuẩn hoá vừa là lợi ích vừa là chi phí.}
So với \texttt{none}, \texttt{standard} tăng $\mathrm{ARI}^{\text{toàn}}$ gấp
gần ba lần ($0.049 \to 0.146$) nhưng làm thông lượng giảm $39$ lần
($39\,860 \to 1\,022$ mẫu/s). Nguyên nhân nằm ở cột cuối: chuẩn hoá trải dữ liệu
ra nên số micro-cluster trung bình tăng từ $4.9$ lên $109.5$, và chi phí mỗi mẫu
của DBSTREAM tăng theo $|\mathcal{M}|$. Hình~\ref{fig:ablation} trình bày lại
hai cột chất lượng của Bảng~\ref{tab:ablation} cạnh nhau.

\begin{figure}[htbp]
\centering
\begin{tikzpicture}
\begin{axis}[
  width=14cm, height=6.4cm,
  ybar=1.2pt,
  bar width=8pt,
  symbolic x coords={standard,rolling,none,robust,adaptive,minmax},
  xtick=data,
  x tick label style={font=\scriptsize,rotate=18,anchor=east},
  ylabel={ARI},
  ymin=-0.15, ymax=0.62,
  legend pos=north west,
  legend cell align=left,
  legend style={font=\small},
  grid=both, grid style={gray!18},
  label style={font=\small},
  tick label style={font=\small},
  extra y ticks={0},
  extra y tick labels={},
  extra y tick style={grid=major, grid style={black!55,thick,dashed}},
]
\addplot[fill=blue!55,draw=blue!70!black] coordinates {
(standard,0.1458)(rolling,0.1225)(none,0.0491)(robust,0.0039)(adaptive,-0.0498)(minmax,-0.0850)
};
\addlegendentry{ARI toàn stream (đáng tin)}
\addplot[fill=orange!65,draw=orange!80!black] coordinates {
(standard,0.0228)(rolling,0.0182)(none,0.0216)(robust,0.0054)(adaptive,0.0608)(minmax,0.5377)
};
\addlegendentry{ARI trung bình per-window (thoái hoá)}
\end{axis}
\end{tikzpicture}
\caption{Cùng sáu lần chạy, hai cách chấm điểm, hai bảng xếp hạng đối nghịch.
Cột cam của \texttt{minmax} vọt lên $0.5377$ hoàn toàn nhờ các cửa sổ thoái hoá;
cột xanh cho thấy phân hoạch thật của nó còn tệ hơn ngẫu nhiên.}
\label{fig:ablation}
\end{figure}

\subsection{Hệ quả trên một lượt chạy đầy đủ: hai bảng xếp hạng ngược nhau}
\label{sec:fullbench}

Đây là kết quả tổng hợp (Bảng~\ref{tab:fullperf} cho hiệu năng,
Bảng~\ref{tab:fullquality} cho chất lượng), và cũng là minh chứng mạnh nhất cho
Mệnh đề~\ref{prop:degenerate}. Bốn thuật toán chạy trên cùng $50\,000$ mẫu đầu
KDDCup99, cùng \texttt{standard} scaler, $m=1000$, $25$ checkpoint, seed
$\{1,2,3\}$ cho thuật toán ngẫu nhiên.

\begin{table}[htbp]
\centering
\small
\caption{Hiệu năng và độ phức tạp mô hình. Cột ``Chạy'' là số lần chạy thật:
DenStream và DBSTREAM tất định nên chỉ chạy một seed, và độ lệch chuẩn $0$ của
chúng là hệ quả định nghĩa chứ không phải bằng chứng về độ ổn định.}
\label{tab:fullperf}
\begin{tabular}{@{}lrrrrrr@{}}
\toprule
\textbf{Thuật toán} & \textbf{Chạy} & \textbf{Runtime (s)} & \textbf{mẫu/s}
 & \textbf{Model MB} & $\overline{|\mathcal{M}|}$ & \textbf{Macro} $K$ \\
\midrule
CluStream    & 3 & $162.7 \pm 5.1$ & $308 \pm 10$    & 0.913 & 100.0 & 5 \\
DenStream    & 1 & $2.5$           & $19\,921$       & 0.012 & $\mathbf{1.0}$ & $\mathbf{1}$ \\
DBSTREAM     & 1 & $53.2$          & $941$           & 0.657 & 109.5 & 1 \\
STREAMKMeans & 3 & $5.6 \pm 0.1$   & $8\,928 \pm 88$ & 3.246 & --- & 10 \\
\bottomrule
\end{tabular}
\end{table}

\begin{table}[htbp]
\centering
\small
\caption{Chất lượng theo hai cách chấm. Thứ hạng trong ngoặc. Với
$\rho_{\text{deg}}=72\%$ trên luồng này, cột trái là cột đáng tin.}
\label{tab:fullquality}
\begin{tabular}{@{}lrrrr@{}}
\toprule
\multirow{2}{*}{\textbf{Thuật toán}}
 & \multicolumn{2}{c}{\textbf{Toàn stream (đáng tin)}}
 & \multicolumn{2}{c}{\textbf{Per-window (thoái hoá)}} \\
\cmidrule(lr){2-3}\cmidrule(lr){4-5}
 & ARI & AMI & mean ARI & median ARI \\
\midrule
DBSTREAM     & $\mathbf{0.1458}$~(1) & 0.1530 & 0.0228 \;(4) & 0.0000 \\
STREAMKMeans & $0.0734 \pm 0.0719$~(2) & 0.2245 & $0.3488 \pm 0.0076$~(2) & 0.0001 \\
CluStream    & $0.0282 \pm 0.0092$~(3) & 0.1298 & $0.1533 \pm 0.0466$~(3) & 0.0000 \\
DenStream    & $\mathbf{0.0000}$~(4) & 0.0000 & $\mathbf{0.7200}$~(1) & 1.0000 \\
\bottomrule
\end{tabular}
\end{table}

\paragraph{Đảo ngược hoàn toàn ở hai đầu bảng.}
DenStream ở cấu hình này \emph{đã sụp hoàn toàn}: đúng một micro-cluster, đúng
một macro-cluster, ARI toàn stream bằng chính xác $0.0000$, AMI bằng $0.0000$.
Về mọi mặt, nó không phân cụm gì cả. Vậy mà nó đứng \textbf{nhất} về ARI
per-window với $0.7200$ --- gấp $32$ lần DBSTREAM, mô hình thực sự tốt nhất
theo chỉ số toàn stream. Trung vị ARI per-window của nó bằng $1.0000$.

Điều này không phải là một sự trùng hợp mà là hệ quả trực tiếp của
Mệnh đề~\ref{prop:degenerate}: một mô hình gán \emph{mọi} điểm vào \emph{một}
cụm sẽ đạt điểm tối đa trên \emph{mọi} cửa sổ thoái hoá, và $72\%$ số cửa sổ ở
đây thoái hoá. DenStream là mô hình under-segment cực đại, nên nó là mô hình
thắng cực đại theo chỉ số đã bị thoái hoá.

Tương quan hạng Spearman giữa hai cột là
\begin{equation}
  \rho_s \;=\; 1 - \frac{6\sum_i d_i^2}{K(K^2-1)}
  \;=\; 1 - \frac{6 \cdot 18}{4 \cdot 15} \;=\; -0.8,
\end{equation}
tức hai chỉ số gần như phản bác nhau. Đây chính là lý do bảng chẩn đoán cửa sổ
được in \emph{trước} bảng time-aware trong báo cáo tự động, kèm cảnh báo khi
$\rho_{\text{deg}} \ge 0.3$.

\paragraph{Đọc bảng hiệu năng cho đúng.}
Bảng~\ref{tab:fullperf} cũng nhắc rằng thông lượng cao không phải lúc nào cũng
là tin tốt: DenStream đạt $19\,921$ mẫu/s \emph{vì} nó chỉ duy trì một
micro-cluster --- tốc độ ở đây là triệu chứng của sự sụp đổ, không phải của
hiệu quả. Ngược lại, CluStream chậm nhất ($308$ mẫu/s) vì nó bị ràng buộc luôn
giữ đủ $100$ micro-cluster (\texttt{max\_micro\_clusters}) và tái phân cụm định
kỳ. Cột $\overline{|\mathcal{M}|}$ là biến giải thích cho cả hai, và
Hình~\ref{fig:scalability} cho thấy quan hệ đó dọc theo luồng.

\begin{figure}[htbp]
\centering
\begin{tikzpicture}
\begin{axis}[
  width=7.1cm, height=5.4cm,
  xlabel={số mẫu đã xử lý},
  ylabel={runtime luỹ kế (s)},
  xmin=0, xmax=50000,
  ymin=0,
  legend pos=north west,
  legend style={font=\scriptsize},
  legend cell align=left,
  grid=both, grid style={gray!18},
  label style={font=\small}, tick label style={font=\scriptsize},
  scaled x ticks=false,
  xticklabel style={/pgf/number format/fixed,/pgf/number format/1000 sep={\,}},
]
\addplot[thick,red!70!black,mark=*,mark size=1pt] coordinates {
(2000,7.62)(8000,31.40)(14000,48.69)(20000,70.88)(26000,92.91)(32000,114.59)
(38000,136.84)(44000,155.18)(50000,168.36)};
\addlegendentry{CluStream}
\addplot[thick,blue!70!black,mark=square*,mark size=1pt] coordinates {
(2000,0.92)(8000,6.59)(14000,11.87)(20000,17.69)(26000,24.92)(32000,31.78)
(38000,40.06)(44000,49.74)(50000,53.15)};
\addlegendentry{DBSTREAM}
\addplot[thick,green!45!black,mark=triangle*,mark size=1.2pt] coordinates {
(2000,0.22)(8000,0.91)(14000,1.57)(20000,2.25)(26000,2.91)(32000,3.60)
(38000,4.27)(44000,4.97)(50000,5.65)};
\addlegendentry{STREAMKMeans}
\addplot[thick,orange!85!black,mark=diamond*,mark size=1.2pt] coordinates {
(2000,0.08)(8000,0.39)(14000,0.69)(20000,0.99)(26000,1.28)(32000,1.58)
(38000,1.88)(44000,2.19)(50000,2.51)};
\addlegendentry{DenStream}
\end{axis}
\end{tikzpicture}
\hfill
\begin{tikzpicture}
\begin{axis}[
  width=7.1cm, height=5.4cm,
  xlabel={số mẫu đã xử lý},
  ylabel={số micro-cluster $|\mathcal{M}|$},
  xmin=0, xmax=50000,
  ymin=0, ymax=240,
  legend pos=north west,
  legend style={font=\scriptsize},
  legend cell align=left,
  grid=both, grid style={gray!18},
  label style={font=\small}, tick label style={font=\scriptsize},
  scaled x ticks=false,
  xticklabel style={/pgf/number format/fixed,/pgf/number format/1000 sep={\,}},
]
\addplot[thick,red!70!black,mark=*,mark size=1pt] coordinates {
(2000,100)(8000,100)(14000,100)(20000,100)(26000,100)(32000,100)
(38000,100)(44000,100)(50000,100)};
\addlegendentry{CluStream (bị chặn trên)}
\addplot[thick,blue!70!black,mark=square*,mark size=1pt] coordinates {
(2000,91)(8000,130)(14000,75)(20000,132)(26000,103)(32000,123)
(38000,117)(44000,206)(50000,2)};
\addlegendentry{DBSTREAM (tăng rồi tỉa)}
\addplot[thick,orange!85!black,mark=diamond*,mark size=1.2pt] coordinates {
(2000,1)(8000,1)(14000,1)(20000,1)(26000,1)(32000,1)
(38000,1)(44000,1)(50000,1)};
\addlegendentry{DenStream (đã sụp)}
\end{axis}
\end{tikzpicture}
\caption{Trái: runtime luỹ kế tăng \emph{tuyến tính} theo số mẫu ở cả bốn thuật
toán --- không có thuật toán nào suy biến về bậc cao hơn trên khoảng này. Phải:
số micro-cluster giải thích độ dốc của đồ thị trái. CluStream bị chặn cứng ở
$100$; DBSTREAM dao động giữa $75$ và $206$ theo nhịp phai--tỉa và tụt xuống $2$
ở cuối; DenStream nằm bẹp ở $1$ suốt luồng. Quỹ đạo dao động của DBSTREAM cũng
minh hoạ Nhận xét~\ref{rem:memcost}: với mô hình có tỉa, đỉnh đo tại checkpoint
là một chặn dưới của đỉnh thật.}
\label{fig:scalability}
\end{figure}

\subsection{Framework phơi bày quan hệ tham số--thang đo}
\label{sec:eps}

Chạy DenStream với $\varepsilon = 16.0$ trên KDDCup99 đã chuẩn hoá
(Nhận xét~\ref{rem:eps}) cho kết quả đúng như dự đoán lý thuyết
(Bảng~\ref{tab:eps}):

\begin{table}[htbp]
\centering
\small
\caption{DenStream, $\varepsilon=16.0$, \texttt{standard}, KDDCup99, $n=5000$.}
\label{tab:eps}
\begin{tabular}{@{}lrrrr@{}}
\toprule
\textbf{Thuật toán} & $\max|\mathcal{M}|$ & \textbf{Macro} $K$ & \textbf{ARI} & \textbf{Purity} \\
\midrule
DenStream ($\varepsilon=16$) & $\mathbf{1}$ & 1 & $\mathbf{0.0000}$ & 0.9992 \\
DBSTREAM  & 132 & 52 & 0.0002 & 0.9992 \\
STREAMKMeans & --- & 10 & 0.0016 & 0.9992 \\
\bottomrule
\end{tabular}
\end{table}

Toàn bộ luồng sụp về đúng một micro-cluster và ARI bằng $0$. Đáng chú ý là
purity vẫn báo $0.9992$ cho \emph{cả ba} thuật toán, kể cả cái đã sụp hoàn toàn
--- một minh hoạ cụ thể cho cảnh báo ở~\eqref{eq:purity} rằng purity không hiệu
chỉnh theo may rủi và vô dụng khi phân bố lớp cực kỳ mất cân bằng.

% =====================================================================
\section{Hai lỗi của River: ghi lại, không vá}
\label{sec:riverbugs}
% =====================================================================

Trong quá trình xây dựng framework, hai lỗi của River lộ ra. Cây nguồn River
\textbf{không bị chỉnh sửa}: nguyên tắc của dự án là framework đo thư viện
\emph{thật}, chứ không phải một bản đã vá cục bộ --- một bản vá cục bộ sẽ biến
mất ở lần cập nhật hoặc revert kế tiếp, và điều đó đã thực sự xảy ra một lần.
Hai lỗi được ghi trong \texttt{RIVER\_BUGS.md} kèm cách tái hiện tối thiểu, và
được né ở phía framework.

\subsection{BUG-1 --- StandardScaler cửa sổ trượt sinh ra số phức}
\label{sec:bug1}

\texttt{StandardScaler(window\_size=10000)} làm DBSTREAM dừng với
\texttt{TypeError: VectorDict values must be real numbers}, báo lỗi ở tận hàm
khoảng cách Minkowski. Nguyên nhân là điều kiện bảo vệ trong~\eqref{eq:zscore}
được cài dưới dạng kiểm tra ``khác không'' thay vì ``lớn hơn không''. Trên đặc
trưng \texttt{num\_shells} --- hằng số $0$ trong suốt cửa sổ --- phương sai trượt
trôi xuống một giá trị \emph{âm} cực nhỏ do triệt tiêu dấu phẩy động:
\[
  \sigma^2 = -2.0095 \times 10^{-18}
  \;\Longrightarrow\;
  (\sigma^2)^{1/2} = 1.4176\times10^{-9}\,i
  \;\Longrightarrow\;
  z = 2.517\times10^{-26} - 4.111\times10^{-10}\,i .
\]
Trong Python, luỹ thừa $0.5$ của một float âm trả về \emph{số phức} chứ không
raise, nên giá trị này đi thẳng vào mô hình và chỉ bộc lộ rất xa nguồn gốc.

\textbf{Sửa ở thượng nguồn:} đổi điều kiện thành $\sigma^2 > 0$ ở cả hai nhánh.

\textbf{Né ở framework:} bước \texttt{standard\_safe} bọc scaler của River và ép
mọi giá trị không phải số thực hữu hạn về $0.0$ --- đúng ý định ban đầu của điều
kiện bảo vệ với một đặc trưng không có độ tán. Bước \texttt{standard} vẫn giữ
nguyên hành vi River để hai bên so sánh được. Đo trên $20\,000$ mẫu KDDCup99:

\begin{center}
\small
\begin{tabular}{@{}lr@{}}
\toprule
\textbf{Bước} & \textbf{Số giá trị phức} \\
\midrule
\texttt{standard(window\_size=10000)} & $5\,887$ \\
\texttt{standard\_safe(window\_size=10000)} & $\mathbf{0}$ \\
\bottomrule
\end{tabular}
\end{center}

\subsection{BUG-2 --- \texorpdfstring{\texttt{seed}}{seed} của KMeans không đủ để tái lập}
\label{sec:bug2}

Cùng \texttt{seed}, cùng dữ liệu, cùng mã --- nhưng kết quả khác nhau giữa các
tiến trình. Đo qua chính framework này trên luồng tổng hợp
($n=3000$, $d=10$, $k=5$, STREAMKMeans, \texttt{seed=1}):

\begin{center}
\small
\begin{tabular}{@{}lr@{}}
\toprule
\texttt{PYTHONHASHSEED} & \textbf{ARI} \\
\midrule
0 & 0.2360384106 \\
1 & 0.5360306622 \\
7 & 0.5784372243 \\
\bottomrule
\end{tabular}
\end{center}

Luồng dữ liệu giống hệt nhau ở cả ba lần --- cùng nhãn, cùng giá trị sau tiền xử
lý. Chỉ đầu ra của mô hình khác. Nguyên nhân là sự kết hợp của hai chi tiết
trong \texttt{predict\_one}:

\begin{enumerate}
  \item khoảng cách được tính bằng cách duyệt hợp các key dưới dạng
        \emph{set literal}, mà thứ tự duyệt set của khoá chuỗi phụ thuộc vào
        \texttt{PYTHONHASHSEED} --- giá trị Python ngẫu nhiên hoá theo tiến trình;
  \item các tâm cụm được giữ trong \texttt{defaultdict(rand\_gauss)}, nên
        \emph{chạm} vào một key mà tâm chưa thấy sẽ \textbf{rút giá trị kế tiếp
        từ RNG đã seed} --- đây chính là cơ chế khởi tạo tâm.
\end{enumerate}

Thứ tự duyệt vì thế quyết định \emph{toạ độ ngẫu nhiên nào rơi vào đặc trưng
nào}. Tâm khởi tạo khác nhau, và trong k-means một khác biệt sớm trong phép gán
lan thành một phân cụm hoàn toàn khác. STREAMKMeans thừa hưởng lỗi này vì nó tạo
một \texttt{KMeans} mới cho mỗi khối.

\textbf{Sửa ở thượng nguồn:} duyệt hợp các key theo thứ tự xác định, vật chất hoá
key của tâm trước khi chạm vào key mới:
\begin{lstlisting}[language=Python]
keys = list(center)
keys += [k for k in x if k not in center]
return sum((abs(center[k] - x.get(k, 0))) ** self.p for k in keys)
\end{lstlisting}

\textbf{Né ở framework:} không có cách sửa từ bên ngoài, nên framework làm cho
lỗi \emph{hiện ra} thay vì che đi. \texttt{capture\_environment()} ghi lại
\texttt{PYTHONHASHSEED} vào mọi báo cáo, và runner cảnh báo khi giá trị này không
được ghim. Bảo đảm mà framework có thể đứng sau là bảo đảm yếu hơn nhưng đúng:
\emph{ghim hash seed thì kết quả tái lập được}.

\begin{lstlisting}[language=bash]
PYTHONHASHSEED=0 python run_benchmark.py --dataset insects --seeds 1,2,3
\end{lstlisting}

Một bài kiểm thử khẳng định trường hợp đã ghim là tái lập, và một bài khác khẳng
định \emph{sự nhạy cảm vẫn còn} --- nếu bài thứ hai bắt đầu fail thì River đã sửa
lỗi và mục này nên được đóng lại.

\begin{remark}[Vì sao điều này quan trọng với một benchmark]
Tham số \texttt{seed} quảng cáo một tính tái lập mà nó không cung cấp giữa các
tiến trình. Chạy lặp nhiều seed rồi báo mean $\pm$ std sẽ trộn phương sai do seed
với phương sai do thứ tự hash, và kết quả không tái lập được trên máy khác hay
ngày khác nếu không biết hash seed đang có hiệu lực. Đây đúng là loại sai lệch mà
bản ghi môi trường (Bảng~\ref{tab:env}) sinh ra để phát hiện.
\end{remark}

% =====================================================================
\section{Cách chạy}
\label{sec:usage}
% =====================================================================

Một entry point duy nhất, \texttt{run\_benchmark.py}, điều khiển hoàn toàn qua
dòng lệnh. Framework \emph{không} chọn siêu tham số hay dataset thay người dùng:
mọi giá trị mặc định chỉ là chỗ khởi động chạy được, và tham số không được nêu
trên dòng lệnh sẽ rơi về mặc định của chính River.

\begin{lstlisting}[language=bash]
# defaults: synthetic stream, sample-partitioning algorithms, no preprocessing
python run_benchmark.py

# your parameters, as JSON
python run_benchmark.py --dataset kdd99 --n-samples 50000 \
    --algorithms DenStream,DBSTREAM \
    --params '{"DenStream": {"epsilon": 0.5, "mu": 2.5}}' \
    --scalers standard

# compare preprocessing chains for one algorithm
python run_benchmark.py --algorithms DBSTREAM \
    --scalers 'standard,standard_safe,variance_threshold|standard'

# a text stream and the clusterer that consumes token counts
python run_benchmark.py --dataset sms_spam --algorithms TextClust \
    --scalers bag_of_words --scaler-params '{"bag_of_words": {"on": "body"}}'

# reproducible figures -- see the note on PYTHONHASHSEED
PYTHONHASHSEED=0 python run_benchmark.py --dataset insects --seeds 1,2,3
\end{lstlisting}

Kết quả ghi ra \texttt{benchmark\_out/<dataset>/}: \texttt{report.md} (bảng +
bản ghi môi trường), \texttt{results.json} (toàn bộ quỹ đạo cửa sổ và
checkpoint), và các đồ thị PNG.

Dùng bằng mã, khi cần điều khiển chi tiết hơn dòng lệnh cho phép:

\begin{lstlisting}[language=Python]
from benchmark.experiments import BenchmarkConfig, Setting, run_setting

setting = Setting("DBSTREAM",
                  params={"clustering_threshold": 1.0},
                  scaler="impute|one_hot_categorical|standard",
                  scaler_params={"standard": {"window_size": 10_000}})
result = run_setting(setting, dataset,
                     BenchmarkConfig(window_size=1000, n_checkpoints=25), seed=1)
\end{lstlisting}

\noindent
\texttt{STRUCTURE.md} mô tả chức năng từng mô-đun và cách thêm dataset riêng;
\texttt{AGENTS.md} ghi các nguyên tắc làm việc trong repo, trong đó có nguyên
tắc không chỉnh sửa cây nguồn River.

% =====================================================================
\section{Hạn chế và hướng mở rộng}
\label{sec:limits}
% =====================================================================

Ranh giới của framework nằm ở đâu cũng đáng nói rõ. Việc chọn siêu tham số và
cấu hình dataset \emph{không} thuộc phạm vi này: framework phơi bày các núm điều
khiển và các chẩn đoán để người dùng đưa ra lựa chọn đó, chứ không đưa hộ.
Ví dụ ở Mục~\ref{sec:eps} minh hoạ đúng vai trò đó --- nó cho thấy $\varepsilon$
gắn với thang đo sau tiền xử lý và cột $\max|\mathcal{M}|$ phát hiện ngay khi
quan hệ đó bị phá vỡ; nó không đề xuất một giá trị $\varepsilon$.

\begin{enumerate}
  \item \textbf{Hai lỗi của River vẫn còn nguyên.} Chúng được ghi lại chứ không
        được vá (Mục~\ref{sec:riverbugs}). Hệ quả thực tế: dùng
        \texttt{standard\_safe} thay cho \texttt{standard} khi bật cửa sổ trượt
        (Mục~\ref{sec:bug1}), và ghim \texttt{PYTHONHASHSEED} cho mọi con số
        định công bố (Mục~\ref{sec:bug2}).

  \item \textbf{Kích thước cửa sổ so với tự tương quan của luồng.} Kết quả ở
        Mục~\ref{sec:degenerate} cho thấy $m=1000$ quá nhỏ so với độ dài khối
        cùng lớp của KDDCup99. Hai hướng khắc phục: tăng $m$ cho tới khi
        $\rho_{\text{deg}}$ xuống dưới một ngưỡng chấp nhận được, hoặc báo cáo
        chỉ số per-window chỉ trên tập cửa sổ \emph{không} thoái hoá. Hướng thứ
        hai giữ được ý nghĩa time-aware nhưng làm giảm cỡ mẫu hiệu dụng.

  \item \textbf{Đỉnh bộ nhớ là chặn dưới.} Vì phép đo chỉ chạy tại checkpoint
        (Nhận xét~\ref{rem:memcost}), các đỉnh ngắn giữa hai checkpoint bị bỏ
        sót. Với mô hình có tỉa, nên tăng số checkpoint.

  \item \textbf{Khảo sát khả năng mở rộng.} Framework có sẵn hai phép quét ---
        theo số mẫu và theo số chiều --- nhưng báo cáo này chưa trình bày kết
        quả của chúng.

  \item \textbf{Chỉ số nội.} Silhouette streaming chỉ chạy với thuật toán có
        expose tâm cụm. Với DenStream ở cấu hình hiện tại (một cụm duy nhất) nó
        không xác định.
\end{enumerate}

% =====================================================================
\appendix
\section{Bảng ký hiệu}
% =====================================================================

\begin{table}[htbp]
\centering
\small
\begin{tabular}{@{}ll@{}}
\toprule
\textbf{Ký hiệu} & \textbf{Ý nghĩa} \\
\midrule
$\mathcal{S} = ((\xt_t,y_t))_{t=1}^{T}$ & luồng dữ liệu, không xáo trộn \\
$d$ & số chiều đặc trưng \\
$\tilde\xt_t$ & mẫu sau chuẩn hoá (không gian mà mô hình làm việc) \\
$\hat y_t$ & nhãn cụm dự đoán tại bước $t$ \\
$\theta_t$ & trạng thái mô hình sau khi học mẫu $t$ \\
$m$ & kích thước cửa sổ đánh giá (mặc định $1000$) \\
$K$ & số cửa sổ, $K=\lceil n/m \rceil$ \\
$W_k$ & tập chỉ số của cửa sổ thứ $k$ \\
$M^{\mathrm{cum}}_k,\ M^{\mathrm{win}}_k$ & giá trị luỹ kế / theo cửa sổ của độ đo $M$ \\
$n_{ij},\ a_i,\ b_j$ & bảng liên hợp và các biên của nó \\
$\mathcal{M}_t$ & tập micro-cluster tại bước $t$ \\
$\rho_{\text{deg}}$ & tỉ lệ cửa sổ thoái hoá \\
$\tilde p$ & trung vị tỉ lệ lớp áp đảo trong cửa sổ \\
$\lambda,\ \varepsilon,\ \mu,\ \beta$ & tham số DenStream \\
$\alpha$ & hệ số phai của bộ chuẩn hoá thích nghi \\
$W$ & kích thước cửa sổ của scaler trượt \\
\bottomrule
\end{tabular}
\end{table}

\begin{thebibliography}{9}
\bibitem{clustream}
C.~C. Aggarwal, J.~Han, J.~Wang, P.~S. Yu.
\emph{A Framework for Clustering Evolving Data Streams}. VLDB, 2003.

\bibitem{denstream}
F.~Cao, M.~Ester, W.~Qian, A.~Zhou.
\emph{Density-Based Clustering over an Evolving Data Stream with Noise}.
SDM, 2006.

\bibitem{dbstream}
M.~Hahsler, M.~Bolaños.
\emph{Clustering Data Streams Based on Shared Density between Micro-Clusters}.
IEEE TKDE, 2016.

\bibitem{hubert}
L.~Hubert, P.~Arabie. \emph{Comparing Partitions}.
Journal of Classification, 1985.

\bibitem{vinh}
N.~X. Vinh, J.~Epps, J.~Bailey.
\emph{Information Theoretic Measures for Clusterings Comparison}. JMLR, 2010.

\bibitem{river}
J.~Montiel et al. \emph{River: Machine Learning for Streaming Data in Python}.
JMLR, 2021.

\bibitem{welford}
B.~P. Welford. \emph{Note on a Method for Calculating Corrected Sums of Squares
and Products}. Technometrics, 1962.
\end{thebibliography}

\end{document}
